\documentclass[11pt]{article}

\usepackage[utf8]{inputenc}
% microtype 의 폰트 확장은 scalable 폰트를 요구한다. 기본 CM 이 비트맵(Type3)으로
% 잡히는 배포판에서는 lmodern 없이는 "auto expansion is only possible with
% scalable fonts" 로 컴파일이 죽는다.
\usepackage{lmodern}
\usepackage[T1]{fontenc}
\usepackage[margin=1in]{geometry}
\usepackage{amsmath,amssymb}
\usepackage{graphicx}
\usepackage{longtable,booktabs}
\usepackage{microtype}
% svgnames 를 줘야 teal 이 정의된다. 없으면 \cite 하나마다
% "Undefined color `teal'" 이 뜬다 (본편 산출물에서 732회).
\usepackage[svgnames]{xcolor}
% hyperref goes last; it makes \cite clickable through to the bibliography.
\usepackage[colorlinks=true,linkcolor=blue,citecolor=teal,urlcolor=magenta]{hyperref}

% pandoc emits \tightlist but expects the preamble to define it.
\providecommand{\tightlist}{%
  \setlength{\itemsep}{0pt}\setlength{\parskip}{0pt}}

% If the compile times out on Overleaf, comment out \tableofcontents below first.
\title{A Survey on the Optimization of Large Language Model-based Agents}
\author{Generated by AutoSurvey}
\date{}

\begin{document}
\maketitle
\tableofcontents
\newpage

\section{Introduction to Large Language Models and Agents}

\subsection{Introduction to Large Language Models}

Large language models (LLMs) have emerged as a transformative force in the field of natural language processing (NLP), demonstrating unprecedented capabilities in understanding and generating human-like language \cite{ref1}. At their core, LLMs are trained on vast amounts of text data, enabling them to learn patterns, relationships, and structures within language \cite{ref2}. This training process allows LLMs to develop a deep understanding of language, which can be leveraged for a wide range of applications, including text generation, language translation, question answering, and text summarization \cite{ref3}.

One of the key capabilities of LLMs is their ability to learn from large datasets and improve their performance over time through self-supervised learning \cite{ref4}. This approach enables LLMs to discover patterns and relationships within the data, which can be used to generate coherent and contextually relevant text \cite{ref5}. However, LLMs also have limitations, including their reliance on high-quality training data, potential biases, and lack of common sense or real-world experience \cite{ref6}. For instance, LLMs may struggle to understand nuances of language, such as sarcasm, irony, or humor, which can lead to misinterpretation or misgeneration of text \cite{ref7}.

Despite these limitations, LLMs have shown remarkable performance in various NLP tasks, including language translation, question answering, and text generation \cite{ref8}. For example, LLMs have been used to generate coherent and contextually relevant text, such as articles, stories, and dialogues \cite{ref9}. Moreover, LLMs have been applied to various domains, including healthcare, finance, and education, where they have shown promising results in tasks such as clinical text analysis, sentiment analysis, and language understanding \cite{ref3}. The capabilities and limitations of LLMs have significant implications for their applications in real-world scenarios, and it is essential to carefully consider and address the potential risks and challenges associated with their use \cite{ref10}.

The development of more efficient and effective LLMs is an active area of research, with a focus on improving their performance, reducing their computational requirements, and increasing their interpretability \cite{ref11}. This has led to the development of new architectures, training methods, and evaluation metrics, which have improved the state-of-the-art in NLP \cite{ref12}. Furthermore, the application of LLMs in multilingual settings has become an increasingly important area of research, with a focus on developing models that can understand and generate text in multiple languages \cite{ref13}. As the field continues to evolve, it is crucial to ensure that LLMs are developed and applied in a responsible and beneficial manner, with consideration for their potential risks and challenges \cite{ref14}.

In conclusion, large language models have emerged as a powerful tool in the field of NLP, with significant capabilities and limitations \cite{ref1}. While they have shown remarkable performance in various NLP tasks, their potential risks and challenges must be carefully considered and addressed \cite{ref15}. As research in this area continues to advance, it is essential to develop more efficient, effective, and interpretable LLMs that can be applied to real-world scenarios, while minimizing their potential risks and challenges \cite{ref11}. Ultimately, the development of LLMs has the potential to transform various aspects of society, from healthcare and education to finance and entertainment, and it is crucial to ensure that these models are developed and applied in a responsible and beneficial manner \cite{ref16}.

\subsection{Definition and Architecture of Large Language Models}

Large language models (LLMs) have revolutionized the field of natural language processing (NLP) with their ability to learn and represent complex patterns in language. As discussed in the previous section, LLMs have demonstrated unprecedented capabilities in understanding and generating human-like language, making them a crucial component in various NLP applications. At their core, LLMs are a type of neural network designed to process and understand human language. In this subsection, we will delve into the definition, architecture, and types of LLMs, including their components and training methods, to provide a comprehensive understanding of these models.

To begin with, LLMs are defined as neural networks that are trained on large amounts of text data to learn the patterns and structures of language \cite{ref1}. These models are typically trained using a self-supervised approach, where the model is tasked with predicting the next word in a sequence of text, given the context of the previous words. This approach allows the model to learn the patterns and relationships between words, as well as the overall structure of language. The self-supervised training method is a key factor in the success of LLMs, as it enables them to learn from large amounts of data without requiring explicit supervision.

The architecture of LLMs is typically based on the transformer model, which is a type of neural network that is designed specifically for sequence-to-sequence tasks, such as machine translation and text summarization. LLMs are typically built on top of the transformer architecture, with the addition of several key components. These include the use of self-attention mechanisms, which allow the model to attend to different parts of the input sequence simultaneously. LLMs also typically use a large number of parameters, often in the hundreds of millions or even billions, which allows them to capture complex patterns in language.

There are several types of LLMs, each with its own strengths and weaknesses. One of the most common types of LLMs is the autoregressive model, which generates text one word at a time, given the context of the previous words. Autoregressive models are typically trained using a masked language modeling objective, where some of the words in the input sequence are randomly replaced with a mask token, and the model is tasked with predicting the original word. Another type of LLM is the sequence-to-sequence model, which generates text in a more traditional sequence-to-sequence manner. Sequence-to-sequence models are typically trained using a combination of the masked language modeling objective and a next sentence prediction objective, where the model is tasked with predicting whether two input sequences are adjacent in the original text.

The training methods used for LLMs are also an important aspect of their architecture. LLMs are typically trained using a combination of the masked language modeling objective and the next sentence prediction objective. The masked language modeling objective involves randomly replacing some of the words in the input sequence with a mask token, and tasking the model with predicting the original word. The next sentence prediction objective involves tasking the model with predicting whether two input sequences are adjacent in the original text. In addition to these two objectives, LLMs are also often trained using a technique called knowledge distillation. Knowledge distillation involves training a smaller model, called the student model, to mimic the behavior of a larger model, called the teacher model.

In recent years, there has been a growing interest in the development of more efficient and scalable LLMs. One approach to achieving this is through the use of model pruning and quantization. Model pruning involves removing redundant or unnecessary weights and connections in the model, while quantization involves reducing the precision of the model's weights and activations. By applying these techniques, it is possible to reduce the computational requirements of LLMs and make them more suitable for deployment on devices with limited resources \cite{ref17}. Another approach to improving the efficiency of LLMs is through the use of parallelization and distributed training. By splitting the model and data across multiple devices, it is possible to train LLMs much more quickly and achieve state-of-the-art results on a wide range of NLP tasks.

As we move forward to explore the applications of LLMs as agents in various domains, it is essential to understand the underlying architecture and training methods of these models. The following section will discuss the applications of LLMs as agents in natural language processing, computer vision, robotics, and other domains, highlighting their potential to solve real-world problems and transform various aspects of society. By leveraging the capabilities of LLMs, we can unlock new possibilities for human-computer interaction, language understanding, and decision-making, and create more intelligent and autonomous systems that can benefit society as a whole.

\subsection{Applications of Large Language Models as Agents}

Large language models have been increasingly applied as agents in various domains, including natural language processing, computer vision, and robotics. As discussed in the previous section, large language models have revolutionized the field of natural language processing with their ability to learn and represent complex patterns in language. The applications of large language models as agents have shown great potential in solving real-world problems, and their impact is expected to grow as the technology continues to evolve.

One of the primary applications of large language models as agents is in natural language processing. These models have been used to generate human-like text, translate languages, and even create chatbots that can converse with humans \cite{ref18}. For instance, \cite{ref19} demonstrates the use of collaborative large language models to improve human-robot interaction. This application is a natural extension of the capabilities of large language models, which have been shown to be effective in a wide range of natural language processing tasks.

Another significant application of large language models as agents is in computer vision. These models have been used to analyze images, detect objects, and even generate images \cite{ref20}. For example, \cite{ref21} proposes a multimodal language model that integrates visual and linguistic information to generate text and images. This application highlights the potential of large language models to be used in multimodal tasks, where they can process and generate multiple forms of data.

In addition to natural language processing and computer vision, large language models have also been applied as agents in robotics. These models have been used to control robots, generate motion plans, and even learn from human feedback \cite{ref22}. For instance, \cite{ref23} proposes a method for grounding language in robotic affordances, which enables robots to understand and execute natural language instructions. This application demonstrates the potential of large language models to be used in human-robot interaction, where they can enable more natural and intuitive communication between humans and robots.

Large language models have also been applied as agents in other domains, such as programming and coding. These models have been used to generate code, debug programs, and even create new programming languages \cite{ref24}. For example, \cite{ref25} proposes a research approach that investigates the inner mechanics of transformer networks to understand the limitations of large language models in multilingual settings. This application highlights the potential of large language models to be used in a wide range of tasks, from natural language processing to programming and coding.

Furthermore, large language models have been applied as agents in scientific research, such as biology and chemistry. These models have been used to analyze large datasets, generate hypotheses, and even predict protein structures \cite{ref26}. For instance, \cite{ref27} discusses the applications of large language models in bioinformatics, including genomics, proteomics, and personalized medicine. This application demonstrates the potential of large language models to be used in scientific research, where they can help analyze large datasets and generate new hypotheses.

The applications of large language models as agents have also been extended to other domains, such as education and tourism. These models have been used to generate educational materials, create interactive learning environments, and even provide personalized recommendations \cite{ref28}. For instance, \cite{ref29} proposes a system for end-user development of robot applications using large language models. This application highlights the potential of large language models to be used in a wide range of tasks, from education to tourism.

In conclusion, large language models have been increasingly applied as agents in various domains, including natural language processing, computer vision, robotics, programming, scientific research, human-robot interaction, education, and tourism. These applications have shown great potential in solving real-world problems, and their impact is expected to grow as the technology continues to evolve. As we move forward, it is essential to address the challenges and limitations associated with the use of large language models as agents, such as the need for large amounts of data, the potential for errors and biases, and the concern about job displacement. By addressing these challenges, we can unlock the full potential of large language models as agents and enable them to have a significant impact on a wide range of applications. The next section will discuss the challenges and limitations of large language models, including issues related to data quality, scalability, interpretability, and robustness, as well as the need for more efficient training methods and better evaluation metrics.

\subsection{Challenges and Limitations of Large Language Models}

The development of large language models has revolutionized the field of natural language processing, enabling machines to process, understand, and generate human-like text with high accuracy. However, despite their impressive performance, large language models are not without their challenges and limitations. As discussed in the previous section, large language models have been increasingly applied as agents in various domains, including natural language processing, computer vision, and robotics. The applications of large language models as agents have shown great potential in solving real-world problems, and their impact is expected to grow as the technology continues to evolve. Nevertheless, to fully realize the potential of large language models, it is essential to address the challenges and limitations associated with their development and deployment.

One of the major challenges facing large language models is the issue of data quality. The training of large language models requires vast amounts of high-quality data, which can be difficult to obtain, especially for low-resource languages \cite{ref30}. Moreover, the data used to train large language models can be biased, noisy, or incomplete, which can negatively impact the performance of the model \cite{ref31}. For instance, a study found that large language models can perpetuate biases and stereotypes present in the training data, which can lead to unfair and discriminatory outcomes \cite{ref32}. Therefore, it is crucial to develop methods for detecting and mitigating biases in large language models, such as data curation and model auditing \cite{ref33}.

Another significant challenge facing large language models is scalability. As the size of the model increases, the computational resources required to train and deploy the model also increase, making it challenging to scale up the model \cite{ref34}. Moreover, large language models can be prone to overfitting, especially when the training data is limited, which can negatively impact the performance of the model \cite{ref35}. To address these challenges, researchers have proposed various techniques, such as model pruning, quantization, and knowledge distillation, which can help reduce the computational resources required to train and deploy large language models \cite{ref17}.

Interpretability is another significant challenge facing large language models. The complex architecture of large language models can make it difficult to understand how the model is making predictions, which can limit the trust and reliability of the model \cite{ref36}. Moreover, the lack of interpretability can make it challenging to identify and mitigate biases in the model, which can negatively impact the performance of the model \cite{ref37}. To address these challenges, researchers have proposed various techniques, such as attention visualization and feature importance, which can help provide insights into how the model is making predictions \cite{ref38}.

Robustness is another significant challenge facing large language models. Large language models can be vulnerable to adversarial attacks, which can negatively impact the performance of the model \cite{ref39}. Moreover, large language models can be prone to hallucinations, which can lead to the generation of nonsensical or irrelevant text \cite{ref6}. To address these challenges, researchers have proposed various techniques, such as adversarial training and robust optimization, which can help improve the robustness of large language models \cite{ref40}.

Finally, the evaluation of large language models is another significant challenge facing the field. The lack of standardized evaluation metrics and benchmarks can make it challenging to compare the performance of different models, which can limit the development of more accurate and reliable models \cite{ref41}. Moreover, the evaluation of large language models can be computationally expensive, which can limit the frequency and scope of evaluation \cite{ref42}. To address these challenges, researchers have proposed various techniques, such as automated evaluation metrics and benchmarking frameworks, which can help improve the efficiency and effectiveness of evaluation \cite{ref43}.

In conclusion, large language models are powerful tools for natural language processing, but they are not without their challenges and limitations. The issues of data quality, scalability, interpretability, and robustness, as well as the need for more efficient training methods and better evaluation metrics, must be addressed to fully realize the potential of large language models. By developing methods for detecting and mitigating biases, improving the scalability and interpretability of large language models, and developing more effective evaluation metrics and benchmarks, researchers can help improve the accuracy, reliability, and trustworthiness of large language models, which can have a significant impact on a wide range of applications, from language translation and text summarization to chatbots and virtual assistants \cite{ref44}. The optimization of large language model-based agents is a crucial field that requires careful consideration of these challenges and limitations, and will be further discussed in the following section, which will explore the applications and future directions of optimized LLM-based agents.

\subsection{Purpose and Scope of the Survey}

The purpose and scope of this survey are to provide a comprehensive overview of the current state of research in optimizing large language model-based agents, highlighting the importance of this field and its potential impact on real-world applications. The emergence of large language models (LLMs) \cite{ref45} has revolutionized the field of artificial intelligence, enabling the development of autonomous agents that can perform complex tasks. However, the optimization of these agents is crucial to unlock their full potential and ensure their effective deployment in various applications.

The optimization of large language model-based agents is a multidisciplinary field that requires expertise in natural language processing, machine learning, and software engineering. Recent advances in LLMs have led to the development of various optimization techniques, including fine-tuning \cite{ref18}, pruning \cite{ref46}, and knowledge distillation \cite{ref47}. These techniques aim to improve the performance, efficiency, and robustness of LLM-based agents, enabling them to tackle complex tasks and adapt to changing environments.

Building on the challenges and limitations of large language models discussed in the previous section, the optimization of LLM-based agents is essential to address issues related to data quality, scalability, interpretability, and robustness. The scope of this survey is to cover the current state of research in optimizing large language model-based agents, including the various techniques, methodologies, and applications. We will discuss the importance of optimizing LLM-based agents, the challenges and limitations associated with their development, and the potential impact on real-world applications.

The optimization of large language model-based agents has numerous applications in various domains, including natural language processing \cite{ref48}, computer vision \cite{ref49}, and robotics \cite{ref50}. These agents can be used to perform tasks such as text generation, language translation, and sentiment analysis, as well as more complex tasks like decision-making and problem-solving.

The importance of optimizing large language model-based agents cannot be overstated, as they have the potential to revolutionize various industries, including healthcare \cite{ref51}, finance \cite{ref52}, and education \cite{ref8}. However, their development and deployment require careful consideration of various factors, including data quality, model complexity, and ethical implications.

This survey aims to provide a comprehensive overview of the current state of research in optimizing large language model-based agents, highlighting the key challenges, limitations, and opportunities. We will discuss the various optimization techniques, including fine-tuning, pruning, and knowledge distillation, and their applications in different domains. We will also provide an overview of the current research landscape, highlighting the key findings, trends, and future directions in this field.

The survey will cover various aspects of optimizing large language model-based agents, including the importance of data quality \cite{ref53}, model complexity \cite{ref2}, and ethical implications \cite{ref54}. We will also discuss the various applications of optimized LLM-based agents, including natural language processing, computer vision, and robotics.

In conclusion, the optimization of large language model-based agents is a crucial field that requires careful consideration of various factors, including data quality, model complexity, and ethical implications. This survey aims to provide a comprehensive overview of the current state of research in this field, highlighting the key challenges, limitations, and opportunities. We hope that this survey will provide a valuable resource for researchers and practitioners working in this field, enabling them to develop more efficient, effective, and robust large language model-based agents that can be deployed in various applications.

The potential impact of optimizing large language model-based agents on real-world applications is significant, and will be further discussed in the following section, which will explore the applications and future directions of optimized LLM-based agents.

\section{Background and Related Work}

\subsection{Model Architectures}

The current state of research in large language model architectures is a rapidly evolving field, with various models being proposed to improve the performance and efficiency of these models. One of the most popular architectures used in large language models is the transformer-based model \cite{ref55}, which has become the standard architecture used in many natural language processing tasks. The transformer architecture is based on self-attention mechanisms, which allow the model to weigh the importance of different input elements relative to each other. This is different from traditional recurrent neural networks (RNNs), which process input sequences sequentially and use recurrent connections to capture long-range dependencies.

In addition to transformer-based models, recurrent neural networks (RNNs) are also widely used in large language models. RNNs are particularly useful for modeling sequential data, such as text or speech, and have been used in many applications, including language modeling, machine translation, and speech recognition \cite{ref56}. However, RNNs can be computationally expensive and difficult to train, especially for long sequences. Other architectures that have been proposed for large language models include the long short-term memory (LSTM) architecture \cite{ref57} and the gated recurrent unit (GRU) architecture \cite{ref58}. These architectures have been shown to be effective in many natural language processing tasks, but they can be computationally expensive and require large amounts of training data.

Recently, there has been a growing interest in developing more efficient and scalable architectures for large language models. One approach that has been proposed is the use of sparse attention mechanisms, which can reduce the computational cost of the self-attention mechanism \cite{ref59}. Another approach is the use of knowledge distillation, which can transfer knowledge from a large pre-trained model to a smaller model, reducing the computational cost and improving the efficiency of the model \cite{ref60}. The transformer-based architecture has also been extended to other domains, including computer vision and speech recognition. For example, the vision transformer (ViT) architecture has been proposed for image classification tasks \cite{ref61}, and the speech transformer architecture has been proposed for speech recognition tasks \cite{ref62}.

Many other models have been proposed for large language models, including the BERT architecture \cite{ref63}, the RoBERTa architecture \cite{ref1}, and the XLNet architecture \cite{ref64}. These models have been shown to be highly effective in many natural language processing tasks and have achieved state-of-the-art results in many benchmarks. The development of large language models has also led to the creation of many pre-trained models, which can be fine-tuned for specific tasks. These pre-trained models have been shown to be highly effective in many natural language processing tasks and have achieved state-of-the-art results in many benchmarks \cite{ref60}.

The future of large language model architectures is likely to involve the development of even more efficient and scalable models, as well as the application of these models to a wider range of tasks and domains. The use of large language models has the potential to revolutionize many areas of natural language processing, including language translation, text generation, and language understanding. However, there are also many challenges associated with the development and deployment of large language models, including the need for large amounts of training data, the risk of overfitting, and the potential for bias and unfairness. To address these challenges, it is essential to develop large language models that are transparent, explainable, and fair. This requires the development of new technologies and methods that can provide insights into the decision-making processes of large language models, and that can ensure that these models are fair and unbiased.

In conclusion, the current state of research in large language model architectures is a rapidly evolving field, with many different architectures being proposed and developed. The transformer-based architecture is one of the most popular architectures used in large language models, but other architectures, such as RNNs and LSTMs, are also widely used. The development of more efficient and scalable architectures, such as sparse attention mechanisms and knowledge distillation, is an active area of research, and many pre-trained models have been created for fine-tuning on specific tasks. As the field continues to evolve, we can expect to see the development of even more advanced large language models, including models that can handle multiple tasks and domains, and models that can learn from limited amounts of data.

\subsection{Training Methods}

The training of large language models is a complex task that requires careful consideration of the training method and technique being used. Various training methods have been employed for large language models, each with its strengths and weaknesses. One of the most popular training methods is masked language modeling \cite{ref65}, which involves randomly masking a portion of the input tokens and then predicting the masked tokens based on the context. This approach has been shown to be effective in learning contextualized representations of words \cite{ref66}. Another commonly used training method is next sentence prediction \cite{ref66}, which involves predicting whether two sentences are adjacent in the original text. This approach helps to learn sentence-level representations and has been shown to be effective in tasks such as natural language inference \cite{ref67}.

In addition to these two methods, other techniques have been proposed to improve the training of large language models. For example, some studies have explored the use of reinforcement learning \cite{ref68} and generative adversarial networks \cite{ref69} to train language models. These approaches have been shown to be effective in improving the performance of language models on certain tasks, but may require significant computational resources and large amounts of training data. The use of multilingual training data \cite{ref70} has also been explored, which involves training a single model on multiple languages simultaneously. This approach can help to improve the performance of the model on low-resource languages by leveraging the similarities between languages.

Furthermore, some researchers have proposed the use of pre-training objectives that are specifically designed to improve the performance of language models on certain tasks \cite{ref71}. For example, some studies have proposed the use of pre-training objectives that are designed to improve the performance of language models on tasks such as question answering \cite{ref67} and text classification \cite{ref72}. The use of self-supervised learning methods \cite{ref73} has also been explored, which involves training a model on a task that is designed to be similar to the task that the model will be used for, but does not require labeled data.

Other techniques, such as transfer learning \cite{ref74} and meta-learning \cite{ref68}, have also been proposed to improve the performance of language models. Transfer learning involves training a model on one task and then fine-tuning it on another task, while meta-learning involves training a model on a set of tasks and then fine-tuning it on a new task. The use of data augmentation techniques \cite{ref75} has also been explored, which involves generating new training data by applying transformations to the existing data, such as paraphrasing or text noising. Additionally, attention mechanisms \cite{ref76} and graph-based methods \cite{ref77} have been proposed to improve the performance of language models.

The use of multimodal training data \cite{ref78} has also been explored, which involves training a model on multiple forms of data, such as text and images, simultaneously. This approach can help to improve the performance of the model on certain tasks, such as visual question answering and image captioning. Knowledge distillation \cite{ref69} has also been proposed, which involves training a smaller model to mimic the behavior of a larger model. This approach can help to improve the performance of the smaller model on certain tasks.

In recent years, there has been a growing interest in the use of distributed training methods \cite{ref73} and specialized hardware \cite{ref79} to speed up the training process. The use of pre-trained models \cite{ref66} and fine-tuning \cite{ref72} has also become more prevalent. As the field continues to evolve, it is likely that new training methods and techniques will be developed to improve the performance of language models on certain tasks. The choice of training method will depend on the specific task and dataset being used, as well as the computational resources available.

Overall, the training of large language models is a complex task that requires careful consideration of the training method and technique being used. By exploring different training methods and techniques, researchers can develop more effective language models that are capable of achieving state-of-the-art performance on a wide range of tasks. The use of a combination of different training methods and techniques can help to improve the performance of language models on certain tasks, and the choice of training method will depend on the specific task and dataset being used, as well as the computational resources available. As the field continues to evolve, it is likely that the training of large language models will become even more efficient and effective, leading to significant advances in natural language processing and other areas of artificial intelligence.

\subsection{Optimization Techniques}

Optimization techniques play a crucial role in improving the performance of large language models (LLMs) while reducing their computational requirements and memory footprint. As discussed in the previous section, the training of large language models is a complex task that requires careful consideration of the training method and technique being used. To further improve the performance and efficiency of LLMs, optimization techniques such as fine-tuning, pruning, quantization, and knowledge distillation can be applied.

Fine-tuning is a widely used optimization technique for LLMs, which involves adjusting the model's parameters to fit a specific task or dataset \cite{ref80}. This technique has been shown to be effective in improving the performance of LLMs on various natural language processing tasks, such as text classification, sentiment analysis, and question answering. However, fine-tuning can be computationally expensive and requires significant amounts of labeled data. To address these challenges, researchers have proposed various fine-tuning strategies, including parameter-efficient fine-tuning methods, such as Low-Rank Adaptation (LoRA) and Half Fine-Tuning \cite{ref81}.

In addition to fine-tuning, pruning is another optimization technique used to reduce the computational requirements and memory footprint of LLMs \cite{ref82}. Pruning involves removing redundant or unnecessary weights and connections in the model, which can lead to significant reductions in computational costs and memory usage. There are various pruning techniques, including structured pruning, unstructured pruning, and hybrid pruning. Structured pruning involves removing entire layers or blocks of weights, while unstructured pruning involves removing individual weights or connections. Hybrid pruning combines both structured and unstructured pruning techniques to achieve better performance and efficiency.

Quantization is a technique used to reduce the precision of the model's weights and activations, which can lead to significant reductions in memory usage and computational costs \cite{ref83}. Quantization involves representing the model's weights and activations using fewer bits, such as 4-bit or 8-bit integers, instead of 32-bit floating-point numbers. This technique can be applied to both the model's weights and activations, and can be used in combination with other optimization techniques, such as pruning and knowledge distillation.

Knowledge distillation is a technique used to transfer knowledge from a large, pre-trained model to a smaller, student model \cite{ref84}. This technique involves training the student model to mimic the behavior of the teacher model, by minimizing the difference between their output distributions. Knowledge distillation can be used to reduce the size of the model, while preserving its performance, and can be applied to various natural language processing tasks, such as text classification, sentiment analysis, and question answering.

Other optimization techniques, such as mask-based pruning \cite{ref85}, regularized learning \cite{ref86}, adaptive pruning and tuning \cite{ref87}, and data multiplexing \cite{ref88}, have also been proposed to improve the performance and efficiency of LLMs. These techniques can be used individually or in combination to achieve better performance and efficiency, and can be applied to various real-world applications, such as text classification, sentiment analysis, and question answering.

The optimization of LLMs is closely related to their evaluation, which is discussed in the following section. The evaluation metrics and benchmarks used to assess the performance of LLMs, such as perplexity, accuracy, and F1-score, can be used to measure the effectiveness of optimization techniques. By applying optimization techniques, such as fine-tuning, pruning, quantization, and knowledge distillation, researchers and developers can improve the performance and efficiency of LLMs, and reduce their computational requirements and memory footprint. This can enable the deployment of LLMs in real-world applications, where resources are limited, and efficiency is critical.

\subsection{Evaluation Metrics and Benchmarks}

Evaluation metrics and benchmarks play a crucial role in assessing the performance of large language models (LLMs). As discussed in the previous section, optimization techniques such as fine-tuning, pruning, quantization, and knowledge distillation can be used to improve the performance and efficiency of LLMs. To measure the effectiveness of these optimization techniques, evaluation metrics and benchmarks are essential. In this subsection, we will introduce the evaluation metrics and benchmarks used to assess the performance of LLMs, including perplexity, accuracy, and F1-score, as well as benchmarks like GLUE and SuperGLUE.

Perplexity is a widely used evaluation metric for LLMs, which measures the model's ability to predict the next word in a sequence \cite{ref89}. A lower perplexity score indicates better performance. However, perplexity has its limitations, as it only measures the model's ability to predict the next word, and does not take into account other aspects of language understanding, such as syntax, semantics, and pragmatics. To address these limitations, other evaluation metrics such as accuracy and F1-score are used, particularly in tasks such as sentiment analysis, question answering, and text classification \cite{ref90}.

Accuracy measures the proportion of correct predictions made by the model, while F1-score is the harmonic mean of precision and recall. These metrics provide a more comprehensive understanding of the model's performance, as they take into account both the precision and recall of the model's predictions. In addition to these evaluation metrics, benchmarks such as GLUE and SuperGLUE have been developed to assess the performance of LLMs on a wide range of natural language understanding tasks \cite{ref91}. GLUE is a multi-task benchmark that includes nine tasks, such as question answering, sentiment analysis, and text classification. SuperGLUE is an extension of GLUE, which includes eight tasks that are more challenging and require more advanced language understanding capabilities.

Other benchmarks, such as CLUE \cite{ref92} and IndicSUPERB \cite{ref93}, have been developed for specific languages and tasks. These benchmarks provide a more nuanced understanding of the model's performance on specific tasks and languages, and can help identify areas for improvement. However, the evaluation metrics and benchmarks used to assess the performance of LLMs have several limitations and challenges. One of the main challenges is the lack of a single, unified evaluation metric that can capture the model's performance on all tasks and languages \cite{ref94}.

Another challenge is the issue of benchmark leakage, where the training data for the model includes examples from the test set, which can result in overestimation of the model's performance \cite{ref95}. This can be addressed by using techniques such as data augmentation, data anonymization, and data splitting, to ensure that the training and test sets are separate and do not overlap. To address the limitations of current evaluation metrics and benchmarks, there has been a growing interest in developing more advanced metrics and benchmarks that can capture the nuances of human language understanding \cite{ref96}. These metrics and benchmarks aim to evaluate the model's ability to understand the meaning and context of language, rather than just its ability to predict the next word or classify text into predefined categories.

Examples of such metrics include the BLEU score, which measures the similarity between the model's output and a reference translation \cite{ref97}, and the F1-score, which measures the harmonic mean of precision and recall \cite{ref90}. Additionally, benchmarks such as SuperGLUE and CLUE have been developed to evaluate the model's ability to understand the meaning and context of language \cite{ref91}. These metrics and benchmarks provide a more comprehensive understanding of the model's performance and can help identify areas for improvement.

In conclusion, the evaluation metrics and benchmarks used to assess the performance of LLMs play a crucial role in driving progress in the field of natural language processing. While there are limitations and challenges to current methods, they have been instrumental in helping researchers and developers identify areas for improvement, compare the performance of different models, and develop more advanced language understanding capabilities. As the field continues to evolve, it is likely that new evaluation metrics and benchmarks will be developed to address the limitations and challenges of current methods, and to capture the nuances of human language understanding. The development of more advanced evaluation metrics and benchmarks will be essential in evaluating the performance of LLMs and driving progress in the field.

\section{Optimization Techniques for Large Language Models}

\subsection{Fine-Tuning Optimization Techniques}

Fine-tuning optimization techniques are crucial for improving the performance of large language models (LLMs) on specific tasks, building on the foundation established by knowledge distillation and reinforcement learning. These techniques enable the adaptation of pre-trained models to downstream tasks with minimal additional training data, making them highly efficient and effective. In this subsection, we will discuss fine-tuning, pruning, and quantization as key optimization techniques for LLMs, and explore how they can be used in conjunction with knowledge distillation and reinforcement learning to achieve optimal results.

Fine-tuning is a widely used technique for adapting pre-trained LLMs to specific tasks, and can be used to refine the knowledge transferred from a teacher model to a student model through knowledge distillation \cite{ref80}. It involves updating the model's parameters to fit the target task, typically using a small amount of task-specific data. Fine-tuning can be applied to the entire model or to specific layers, depending on the task requirements. For example, \cite{ref98} proposes a prefix-tuning method that updates only a small subset of the model's parameters, achieving comparable performance to full fine-tuning while reducing the number of updated parameters.

Pruning is another optimization technique used to reduce the computational costs and memory requirements of LLMs, which can be particularly useful when combined with reinforcement learning to fine-tune models for specific tasks \cite{ref99}. Pruning involves removing redundant or unnecessary weights and connections in the model, resulting in a more efficient and compact model. Pruning can be applied to the entire model or to specific layers, and various pruning techniques have been proposed, including structured pruning and unstructured pruning. For example, \cite{ref85} proposes a mask-based pruning method that achieves layer-wise uniform structures, resulting in improved performance and reduced computational costs.

Quantization is a technique used to reduce the precision of the model's weights and activations, resulting in a more efficient and compact model, which can be used in conjunction with knowledge distillation to transfer knowledge from a large model to a smaller model \cite{ref100}. Quantization can be applied to the entire model or to specific layers, and various quantization techniques have been proposed, including post-training quantization and quantization-aware training. For example, \cite{ref83} proposes a quantization-based fine-tuning method that achieves efficient fine-tuning of LLMs while reducing the computational costs and memory requirements.

In addition to fine-tuning, pruning, and quantization, other optimization techniques have been proposed to improve the performance of LLMs, such as combining knowledge distillation and reinforcement learning \cite{ref101}. \cite{ref102} proposes a sparsity-preserving parameter-efficient fine-tuning method that achieves improved performance and reduced computational costs. The choice of optimization technique depends on the specific task requirements and the characteristics of the pre-trained model, as well as the desired balance between performance and computational efficiency.

Recent studies have also explored the application of optimization techniques to specific tasks, such as natural language processing and computer vision, and have demonstrated the effectiveness of combining knowledge distillation and reinforcement learning with fine-tuning, pruning, and quantization \cite{ref103}. \cite{ref104} proposes a contextual pruning method that achieves efficient large language models for computer vision tasks. These studies highlight the importance of choosing the right optimization technique for the specific task, and demonstrate the potential for optimization techniques to be used in conjunction with knowledge distillation and reinforcement learning to achieve state-of-the-art results.

In conclusion, fine-tuning optimization techniques are crucial for improving the performance of large language models on specific tasks, and can be used in conjunction with knowledge distillation and reinforcement learning to achieve optimal results. Fine-tuning, pruning, and quantization are key optimization techniques that can be used to adapt pre-trained models to downstream tasks, and the choice of technique depends on the specific task requirements and the characteristics of the pre-trained model. By combining these techniques with knowledge distillation and reinforcement learning, developers can create highly efficient and effective large language models that achieve state-of-the-art performance on a wide range of tasks.

\subsection{Knowledge Distillation and Reinforcement Learning}

Knowledge distillation and reinforcement learning are two essential optimization techniques in the field of large language models, building on the foundation of fine-tuning, pruning, and quantization discussed in the previous section. Knowledge distillation is a method that involves training a smaller model, known as the student, to mimic the behavior of a larger model, known as the teacher \cite{ref105}. This technique enables the transfer of knowledge from a large model to a smaller model, resulting in a more efficient model that maintains a significant portion of the original model's performance. On the other hand, reinforcement learning is a technique that involves training a model to make decisions in an environment to maximize a reward signal \cite{ref106}.

The key strengths of knowledge distillation include its ability to reduce the computational requirements of large language models, making them more suitable for deployment on devices with limited resources \cite{ref107}. Additionally, knowledge distillation can be used to improve the performance of large language models by transferring knowledge from a pre-trained model to a smaller model \cite{ref108}. However, knowledge distillation also has its weaknesses, such as the challenge of selecting the right teacher model and student model \cite{ref109}. Furthermore, knowledge distillation can be sensitive to the choice of hyperparameters, such as the learning rate and batch size, which can significantly impact the performance of the student model \cite{ref110}.

Reinforcement learning, on the other hand, has been used to fine-tune large language models for specific tasks, such as text generation and language translation \cite{ref111}. This technique involves training a model to make decisions in an environment to maximize a reward signal, which can be designed to encourage the model to generate text that is coherent, fluent, and relevant to the task at hand \cite{ref112}. The combination of knowledge distillation and reinforcement learning has been shown to be effective in improving the performance of large language models, while also reducing their computational requirements \cite{ref113}. This approach involves using knowledge distillation to transfer knowledge from a large model to a smaller model, and then fine-tuning the smaller model using reinforcement learning \cite{ref114}.

In recent years, there has been a growing interest in exploring the applications of knowledge distillation and reinforcement learning in optimizing large language models \cite{ref115}. The effectiveness of combining these two techniques has been demonstrated in various studies, which have shown that it can lead to improved performance and reduced computational requirements \cite{ref116}. As the field of large language models continues to evolve, it is likely that knowledge distillation and reinforcement learning will play an increasingly important role in optimizing these models for specific tasks and applications \cite{ref117}. The following section will discuss parameter-efficient fine-tuning methods, which can be used in conjunction with knowledge distillation and reinforcement learning to further improve the performance and efficiency of large language models.

\subsection{Parameter-Efficient Fine-Tuning Methods}

Parameter-efficient fine-tuning methods have gained significant attention in recent years due to their ability to reduce computational costs while maintaining performance. As discussed in the previous section, knowledge distillation and reinforcement learning are two optimization techniques that have been used to improve the performance of large language models. However, these techniques can be computationally expensive, and parameter-efficient fine-tuning methods offer a promising solution to reduce these costs. In this subsection, we will delve into parameter-efficient fine-tuning methods, including low-rank adaptation, sparse updates, and bitfit, and their effectiveness in reducing computational costs.

Low-rank adaptation is a popular parameter-efficient fine-tuning method that has been widely adopted in recent years \cite{ref118}. This method involves decomposing the weight updates of a pre-trained model into two low-rank matrices, which reduces the number of trainable parameters and computational costs. Low-rank adaptation has been shown to be effective in a variety of tasks, including natural language processing and computer vision \cite{ref119}. However, it can be challenging to determine the optimal rank for the low-rank matrices, and a higher rank can lead to overfitting \cite{ref120}.

Sparse updates are another parameter-efficient fine-tuning method that involves updating only a subset of the model's parameters \cite{ref121}. This method can be particularly useful for large language models, which often have a large number of parameters. By updating only a subset of the parameters, sparse updates can reduce computational costs and memory usage. However, sparse updates can be challenging to implement, and the choice of which parameters to update can have a significant impact on performance \cite{ref122}.

Bitfit is a parameter-efficient fine-tuning method that involves updating only the bias terms of a pre-trained model \cite{ref123}. This method is particularly useful for large language models, which often have a large number of bias terms. By updating only the bias terms, bitfit can reduce computational costs and memory usage while maintaining performance. However, bitfit can be limited by the fact that it only updates the bias terms, and may not be able to capture complex patterns in the data \cite{ref124}.

In addition to these methods, there are several other parameter-efficient fine-tuning techniques that have been proposed in recent years. For example, \cite{ref125} proposes a method for pruning low-rank adaptation parameters based on their output evaluation. This method can reduce computational costs and memory usage while maintaining performance. \cite{ref126} proposes a method for breaking the low-rank bottleneck in low-rank adaptation optimization. This method can improve the performance of low-rank adaptation by allowing for higher-rank updates.

\cite{ref127} proposes a computation-efficient low-rank adaptation method for language models. This method can reduce computational costs while maintaining performance. \cite{ref128} proposes a parameter-efficient fine-tuning method that uses singular vectors to reduce computational costs. This method can be particularly useful for large language models, which often have a large number of parameters.

\cite{ref129} proposes a quantization-aware parameter-efficient adaptation method for large-scale pre-trained language models. This method can reduce computational costs and memory usage while maintaining performance. \cite{ref130} proposes an extremely parameter-efficient fine-tuning method for extreme multi-profile scenarios. This method can reduce computational costs and memory usage while maintaining performance.

\cite{ref131} proposes an extreme parameter-efficient fine-tuning method that uses vector banks to reduce computational costs and memory usage. This method can be particularly useful for large language models, which often have a large number of parameters. \cite{ref132} proposes a parameter-efficient uncertainty quantification method for low-rank adaptation. This method can reduce computational costs and memory usage while maintaining performance.

\cite{ref133} proposes a technical report on fine-tuned large language models that rival GPT-4. This report demonstrates the effectiveness of parameter-efficient fine-tuning methods in reducing computational costs and memory usage while maintaining performance. \cite{ref134} proposes a partial rotation method that empowers more parameter-efficient low-rank adaptation. This method can reduce computational costs and memory usage while maintaining performance.

\cite{ref135} proposes a geometric integration method for parameter-efficient fine-tuning. This method can reduce computational costs and memory usage while maintaining performance. \cite{ref136} proposes a random subspace adaptation method for efficient fine-tuning. This method can reduce computational costs and memory usage while maintaining performance.

The use of parameter-efficient fine-tuning methods can be particularly useful when combined with other optimization techniques, such as quantization and pruning, which will be discussed in the following section. By reducing the computational costs and memory usage of large language models, these methods can enable the deployment of these models on resource-constrained devices, which is essential for many real-world applications. In conclusion, parameter-efficient fine-tuning methods have been widely adopted in recent years due to their ability to reduce computational costs while maintaining performance. Low-rank adaptation, sparse updates, and bitfit are popular parameter-efficient fine-tuning methods that have been shown to be effective in a variety of tasks. However, these methods can be challenging to implement, and the choice of which parameters to update can have a significant impact on performance. Further research is needed to develop more effective and efficient parameter-efficient fine-tuning methods that can reduce computational costs and memory usage while maintaining performance.

\subsection{Quantization and Pruning Techniques}

Quantization and pruning are two essential techniques used in model compression to reduce the memory footprint and computational requirements of large language models \cite{ref137}. As discussed in the previous section, parameter-efficient fine-tuning methods have been widely adopted to reduce computational costs while maintaining performance. However, these methods can be further enhanced by combining them with quantization and pruning techniques to enable the deployment of large language models on resource-constrained devices.

Post-training quantization is a technique used to reduce the precision of model weights from 32-bit floating-point numbers to lower-bit integers, such as 8-bit or 4-bit integers \cite{ref138}. This technique is particularly useful when the model is already trained and deployed, and there is no need to retrain the model. Post-training quantization can be performed using various methods, including uniform quantization, non-uniform quantization, and learned quantization \cite{ref139}. However, post-training quantization can result in significant accuracy degradation, especially when the number of quantization bits is low \cite{ref140}.

Quantization-aware training is another technique used to reduce the precision of model weights during training \cite{ref141}. This technique involves training the model with lower-bit precision weights and activations, which can help to reduce the memory footprint and computational requirements of the model. Quantization-aware training can be performed using various methods, including uniform quantization, non-uniform quantization, and learned quantization \cite{ref142}. By combining quantization-aware training with parameter-efficient fine-tuning methods, such as low-rank adaptation and sparse updates, it is possible to achieve significant reductions in computational costs and memory usage while maintaining performance.

Structured pruning is a technique used to remove redundant connections and neurons in a neural network \cite{ref143}. This technique involves identifying the most important connections and neurons in the network and removing the less important ones. Structured pruning can be performed using various methods, including magnitude-based pruning, gradient-based pruning, and regularization-based pruning \cite{ref144}. Unstructured pruning is a technique used to remove individual weights and connections in a neural network \cite{ref145}. This technique involves identifying the most important weights and connections in the network and removing the less important ones. Unstructured pruning can be performed using various methods, including magnitude-based pruning, gradient-based pruning, and regularization-based pruning \cite{ref146}.

The combination of quantization and pruning techniques can result in significant model compression and accuracy degradation \cite{ref147}. However, the combination of these techniques can also result in significant improvements in model efficiency and deployment on resource-constrained devices \cite{ref148}. In recent years, various techniques have been proposed to improve the efficiency and accuracy of quantization and pruning techniques \cite{ref149}. These techniques include the use of knowledge distillation, transfer learning, and meta-learning to improve the accuracy of quantized and pruned models \cite{ref150}.

The application of quantization and pruning techniques has been explored in various domains, including natural language processing \cite{ref151}. These techniques have been used to compress large language models, such as BERT and RoBERTa, and have resulted in significant improvements in model efficiency and deployment on resource-constrained devices \cite{ref152}. As the field continues to evolve, it is likely that we will see the development of more efficient and accurate quantization and pruning techniques, as well as the exploration of new domains and applications for these techniques \cite{ref153}. In the next section, we will discuss the application of these techniques in real-world scenarios and the future directions for research in this area.

\section{Agent-Specific Optimization Methods}

\subsection{Reinforcement Learning from Human Feedback}

Reinforcement learning from human feedback (RLHF) has emerged as a pivotal technique in aligning large language models (LLMs) with human preferences, enabling them to generate more accurate and informative responses. This approach involves learning from human feedback, which can be in the form of preferences, ratings, or corrections, to fine-tune the model and improve its performance. In this subsection, we will delve into the applications, strengths, and weaknesses of RLHF, as well as recent advancements and techniques in this field, and explore how RLHF can be used in conjunction with other optimization methods, such as policy gradient optimization, to achieve better performance.

One of the primary applications of RLHF is in natural language processing (NLP) tasks, such as text generation, language translation, and sentiment analysis. For instance, \cite{ref154} demonstrates the effectiveness of RLHF in improving the performance of LLMs on NLP tasks. The authors propose a large-scale, high-quality, and diversified preference dataset, which is used to train various models, including the reward model UltraRM, chat language model UltraLM-13B-PPO, and critique model UltraCM. The experimental results show that these models outperform existing open-source models, achieving top performance across multiple benchmarks. This highlights the potential of RLHF to improve the performance of LLMs in a variety of tasks, and demonstrates the importance of using high-quality and diverse human feedback.

RLHF has also been applied in other domains, such as computer vision and robotics. For example, \cite{ref155} presents a large dataset of human Atari 2600 replays, which can be used to train RL agents using human feedback. The authors demonstrate that RLHF can be used to learn policies that are comparable to human performance, even in complex and high-dimensional state spaces. This demonstrates the versatility of RLHF and its potential to be applied in a wide range of tasks.

Despite its successes, RLHF also has some weaknesses and challenges. One of the main limitations is the need for large amounts of high-quality human feedback, which can be time-consuming and expensive to collect. Additionally, RLHF can be sensitive to the quality of the feedback, and noisy or biased feedback can negatively impact the performance of the model. \cite{ref156} discusses these challenges and provides insights into the current practices and opportunities for future research in RLHF. To address these challenges, researchers have proposed various methods, such as data smoothing and AI feedback, to improve the efficiency and effectiveness of RLHF.

Recent advancements in RLHF have focused on addressing these challenges and improving the efficiency and effectiveness of the approach. For example, \cite{ref157} proposes a novel method called Iterative Data Smoothing (IDS), which updates the model and data simultaneously during each training epoch. The authors demonstrate that IDS can improve the performance of RLHF and reduce overfitting and overoptimization. Another area of research has focused on developing more efficient and scalable methods for collecting and utilizing human feedback. For instance, \cite{ref158} proposes a technique called RLAIF, which uses AI feedback instead of human feedback to reduce the need for human annotation. The authors demonstrate that RLAIF can achieve similar performance to RLHF, while reducing the need for human feedback by up to 90\%.

In addition to these advancements, there has been a growing interest in developing more robust and generalizable RLHF methods. For example, \cite{ref159} proposes a novel approach called MaxMin-RLHF, which learns a mixture of preference distributions via an expectation-maximization algorithm and proposes a MaxMin alignment objective for policy learning. The authors demonstrate that MaxMin-RLHF can improve the performance of RLHF and provide a more equitable alignment of LLMs with diverse human preferences. This highlights the potential of RLHF to be used in a variety of tasks, and demonstrates the importance of developing more robust and generalizable methods.

Overall, RLHF has emerged as a powerful technique for aligning LLMs with human preferences, with applications in NLP, computer vision, robotics, and healthcare. While there are challenges and limitations to this approach, recent advancements have focused on addressing these challenges and improving the efficiency and effectiveness of RLHF. As the field continues to evolve, we can expect to see further innovations and improvements in RLHF, enabling more accurate and informative responses from LLMs. Furthermore, the combination of RLHF with other optimization methods, such as policy gradient optimization, has the potential to achieve even better performance, and is an area of research that is worth exploring in the future. \cite{ref160} discusses the applications of RLHF in healthcare, including dynamic treatment regimes, automated medical diagnosis, and personalized medicine. The authors highlight the potential of RLHF to improve patient outcomes and reduce healthcare costs, while also discussing the challenges and limitations of this approach in healthcare.

In addition, \cite{ref161} provides a comprehensive survey of deep reinforcement learning methods for building human-level agents, including RLHF. The authors discuss the challenges and opportunities of using RLHF to develop more intelligent and autonomous systems, and highlight the need for further research in this area. Moreover, \cite{ref162} provides a comprehensive review of RL algorithms and applications in healthcare and robotics, including RLHF. The authors discuss the strengths and weaknesses of RLHF, as well as its potential applications in these domains, and highlight the need for further research to address the challenges and limitations of this approach.

In conclusion, RLHF has emerged as a powerful technique for aligning LLMs with human preferences, with applications in NLP, computer vision, robotics, and healthcare. While there are challenges and limitations to this approach, recent advancements have focused on addressing these challenges and improving the efficiency and effectiveness of RLHF. As the field continues to evolve, we can expect to see further innovations and improvements in RLHF, enabling more accurate and informative responses from LLMs. \cite{ref163} provides a comprehensive platform and benchmark suite for RLHF, which can be used to evaluate and compare different RLHF methods, and \cite{ref164} provides a systematic review of the current state of knowledge on RL-enhanced LLMs, including RLHF.

\subsection{Policy Gradient Optimization}

Policy gradient optimization methods have been widely used in reinforcement learning to optimize the policies of agents, including large language model-based agents. As discussed in the previous section, reinforcement learning from human feedback (RLHF) has emerged as a pivotal technique in aligning large language models with human preferences. Policy gradient optimization methods can be used in conjunction with RLHF to achieve better performance. The theoretical foundations of policy gradient optimization methods are based on the concept of policy gradients, which represent the gradient of the expected cumulative reward with respect to the policy parameters \cite{ref165}. The policy gradient theorem provides a way to compute the policy gradient, which is essential for optimizing the policy using gradient-based methods \cite{ref166}.

One of the advantages of policy gradient optimization methods is that they can handle high-dimensional action spaces and complex policies \cite{ref167}. This makes them particularly suitable for optimizing large language model-based agents, which often have complex policies and high-dimensional action spaces. Additionally, policy gradient optimization methods can be used with a variety of exploration strategies, such as entropy regularization and curiosity-driven exploration \cite{ref168}. These exploration strategies can be used to improve the efficiency and effectiveness of policy gradient optimization methods, especially when combined with RLHF.

However, policy gradient optimization methods also have some limitations. One of the main limitations is that they can be sensitive to the choice of hyperparameters, such as the learning rate and the entropy regularization coefficient \cite{ref169}. Additionally, policy gradient optimization methods can be computationally expensive, especially when dealing with large language models and complex policies \cite{ref170}. To address these limitations, researchers have proposed various variants of policy gradient optimization methods, including trust region policy optimization (TRPO) and proximal policy optimization (PPO) \cite{ref169}.

Despite these limitations, policy gradient optimization methods have been successfully applied to optimize large language model-based agents in a variety of tasks, including natural language processing, computer vision, and robotics \cite{ref165}. For example, \cite{ref165} used policy gradient optimization to optimize a large language model-based agent for a retrospective language modeling task, achieving state-of-the-art results. Similarly, \cite{ref171} used policy gradient optimization to optimize a large language model-based agent for a financial trading task, achieving significant improvements over baseline methods.

In addition to these applications, policy gradient optimization methods have also been used to optimize large language model-based agents in multi-agent settings \cite{ref172}. For example, \cite{ref172} used policy gradient optimization to optimize a large language model-based agent in a multi-agent setting, achieving significant improvements over baseline methods. These results demonstrate the potential of policy gradient optimization methods to be used in a variety of tasks and settings, and highlight the need for further research in this area.

As we will discuss in the following section, multi-agent collaboration and optimization techniques have gained significant attention in recent years, particularly with the emergence of large language models \cite{ref173}. Policy gradient optimization methods can be used in conjunction with these techniques to achieve better performance and improve the overall effectiveness of large language model-based agents. By combining policy gradient optimization methods with other optimization methods, such as value-based methods and model-based methods, researchers and practitioners can develop more effective methods for optimizing large language model-based agents \cite{ref168}.

In conclusion, policy gradient optimization methods are a powerful tool for optimizing large language model-based agents. While they have some limitations, they have been successfully applied to a variety of tasks and have achieved state-of-the-art results. By understanding the theoretical foundations and advantages of policy gradient optimization methods, as well as their limitations and challenges, researchers and practitioners can develop more effective methods for optimizing large language model-based agents \cite{ref166}. Furthermore, policy gradient optimization methods can be combined with other optimization methods and techniques, such as RLHF and multi-agent collaboration, to achieve better performance and improve the overall effectiveness of large language model-based agents.

\subsection{Multi-Agent Collaboration and Optimization}

Multi-agent collaboration and optimization techniques have gained significant attention in recent years, particularly with the emergence of large language models (LLMs) \cite{ref173}. As we discussed in the previous section, policy gradient optimization methods are a powerful tool for optimizing large language model-based agents. However, these methods can be limited when applied to complex tasks that require the coordination of multiple agents. This is where multi-agent collaboration and optimization techniques come in, enabling multiple agents to work together to achieve a common goal, leveraging their individual strengths and capabilities to improve overall performance.

One of the key benefits of multi-agent collaboration is the ability to tackle complex tasks that are beyond the capabilities of individual agents \cite{ref174}. By working together, agents can share knowledge, expertise, and resources, leading to more effective and efficient problem-solving. For example, in the context of cooperative pathfinding, multiple agents can collaborate to find optimal paths in complex environments \cite{ref175}. This is particularly relevant to large language model-based agents, which can be used to optimize complex tasks such as natural language processing and computer vision.

Another benefit of multi-agent collaboration is the ability to adapt to changing environments and situations \cite{ref176}. In dynamic environments, individual agents may not have the necessary information or capabilities to respond effectively to changing conditions. However, by collaborating with other agents, they can share information and adapt to new situations more effectively. For example, in the context of traffic management, multiple agents can collaborate to optimize traffic flow and reduce congestion \cite{ref177}. This ability to adapt to changing environments is crucial for large language model-based agents, which often operate in complex and dynamic environments.

Despite the benefits of multi-agent collaboration, there are also several challenges that need to be addressed. One of the key challenges is the need for effective communication and coordination between agents \cite{ref178}. In order for agents to work together effectively, they need to be able to communicate and coordinate their actions. This can be particularly challenging in complex environments, where agents may have different goals, capabilities, and priorities. For example, in the context of human-robot collaboration, humans and robots need to be able to communicate effectively in order to work together safely and efficiently \cite{ref179}.

Another challenge in multi-agent collaboration is the need for effective optimization techniques \cite{ref180}. In order to optimize the performance of multiple agents, optimization techniques need to be developed that can handle the complexity and uncertainty of multi-agent systems. For example, in the context of cooperative reinforcement learning, optimization techniques need to be developed that can handle the complexity of multiple agents learning together \cite{ref181}. Recent research has made significant progress in addressing these challenges, and we will discuss some of these developments in the following section.

In terms of future research directions, there are several areas that need to be explored further. One area is the development of more effective communication and coordination mechanisms between agents \cite{ref182}. This could involve the development of new communication protocols, such as those based on natural language processing or computer vision. Another area is the development of more effective optimization techniques, such as those based on distributed optimization or multi-objective optimization \cite{ref183}. Additionally, there is a need for more research on the application of multi-agent collaboration and optimization techniques to real-world problems \cite{ref184}.

In conclusion, multi-agent collaboration and optimization techniques have the potential to revolutionize a range of applications, from cooperative robotics to smart cities. By building on the foundations of policy gradient optimization methods, and addressing the challenges of multi-agent collaboration, we can develop more effective and efficient systems that can tackle complex tasks and problems \cite{ref185}. In the following section, we will explore some of the recent developments in multi-agent collaboration and optimization, and discuss their potential applications to large language model-based agents \cite{ref186}.

\section{Evaluation Metrics and Benchmarks for Agent Optimization}

\subsection{Task-Oriented Evaluation Metrics}

Task-oriented evaluation metrics are a crucial aspect of assessing the performance of large language model-based agents, as they enable the evaluation of an agent's ability to complete specific tasks while considering factors such as efficiency and generalization. These metrics are essential in evaluating the effectiveness of agents in real-world applications, where the ability to complete tasks efficiently and generalize to new situations is vital. In this subsection, we will delve into various task-oriented evaluation metrics, including metrics for task completion, efficiency, and generalization, and provide examples of such metrics used in benchmarks, highlighting their significance in the context of large language model-based agents.

To begin with, task completion metrics are used to assess the ability of an agent to complete a specific task. These metrics can be further divided into two categories: success-based metrics and quality-based metrics. Success-based metrics, such as task success rate \cite{ref187}, measure the percentage of tasks that an agent is able to complete successfully. Quality-based metrics, such as task completion score \cite{ref188}, evaluate the quality of the agent's output, such as the accuracy or fluency of the generated text. These metrics provide valuable insights into an agent's ability to complete tasks and are essential in evaluating the performance of large language model-based agents.

In addition to task completion metrics, efficiency metrics are used to evaluate the amount of resources required by an agent to complete a task. These metrics can include measures such as response time \cite{ref189}, which evaluates the time taken by an agent to respond to a user's input, or computational cost \cite{ref190}, which evaluates the amount of computational resources required by an agent to complete a task. Efficiency metrics are crucial in evaluating the performance of large language model-based agents, as they enable the identification of areas where agents can be optimized to reduce resource consumption.

Generalization metrics, on the other hand, are used to evaluate the ability of an agent to generalize to new situations or tasks. These metrics can include measures such as out-of-distribution detection \cite{ref191}, which evaluates an agent's ability to detect when it is faced with a task or situation that is outside its training distribution, or few-shot learning \cite{ref192}, which evaluates an agent's ability to learn from a few examples and generalize to new tasks. Generalization metrics are essential in evaluating the performance of large language model-based agents, as they enable the assessment of an agent's ability to adapt to new situations and tasks.

Several benchmarks have been developed to evaluate the performance of large language model-based agents, including the BIG-bench benchmark \cite{ref193}, which provides a comprehensive evaluation of language generation models. These benchmarks can be used to evaluate the performance of agents on a wide range of tasks, including natural language processing, computer vision, and robotics, and provide valuable insights into the strengths and weaknesses of different agents and models.

The development of task-oriented evaluation metrics is an active area of research, with many new metrics and benchmarks being proposed in recent years. For example, the paper \cite{ref194} proposes a new metric for evaluating the performance of dialogue systems, while the paper \cite{ref195} proposes a new metric for evaluating the performance of text generation models. These new metrics and benchmarks are essential in advancing the field of large language model-based agents and enabling the development of more effective and efficient models.

In the context of large language model-based agents, task-oriented evaluation metrics can be used to evaluate the performance of agents on a wide range of tasks, including natural language processing, computer vision, and robotics. For example, the paper \cite{ref196} proposes a new benchmark for evaluating the performance of vision-language models, while the paper \cite{ref197} proposes a new benchmark for evaluating the performance of single object tracking models. These benchmarks and metrics are crucial in evaluating the performance of large language model-based agents and enabling the development of more effective and efficient models.

In conclusion, task-oriented evaluation metrics are essential in evaluating the performance of large language model-based agents, as they enable the assessment of an agent's ability to complete specific tasks while considering factors such as efficiency and generalization. The development of new metrics and benchmarks is crucial in advancing the field of large language model-based agents and enabling the development of more effective and efficient models. By using these metrics and benchmarks, researchers and developers can gain valuable insights into the strengths and weaknesses of different agents and models, and develop more effective and efficient models that can be used in real-world applications.

\subsection{Benchmarking Frameworks for Large Language Models}

Benchmarking frameworks are a crucial component in the evaluation and optimization of large language models (LLMs), as they provide a standardized way to assess the capabilities of these models and identify areas for improvement. The importance of benchmarking frameworks is closely tied to the task-oriented evaluation metrics discussed in the previous section, as they enable researchers to evaluate the performance of LLMs on various tasks and compare different models. In this subsection, we will delve into the specifics of benchmarking frameworks designed for LLMs, their features, and advantages, highlighting their role in optimizing the performance of these models.

One of the key challenges in benchmarking LLMs is the diversity of tasks they can perform, including natural language processing, computer vision, and multimodal tasks. To address this challenge, benchmarking frameworks like \cite{ref198} have been proposed, providing a comprehensive evaluation suite for LLMs that includes metrics for sequential decoding, parallelization, and setup efficiency. This framework enables researchers to evaluate the performance of LLMs on various tasks and identify areas for optimization, which is essential for developing more efficient and effective models.

The evaluation of LLMs on multilingual tasks is another critical aspect of benchmarking frameworks. \cite{ref199} is a benchmarking framework that provides a comprehensive evaluation of LLMs on multilingual tasks, including translation, question answering, and text classification. This framework enables researchers to evaluate the performance of LLMs on multilingual tasks and identify areas for improvement, which is vital for developing models that can be used in a wide range of languages and applications.

In addition to multilingual tasks, benchmarking frameworks also need to evaluate the performance of LLMs on multimodal tasks, such as vision-language tasks. \cite{ref200} is a benchmarking framework that evaluates the performance of LLMs on multimodal tasks, including visual question answering, image-text retrieval, and visual reasoning. This framework enables researchers to evaluate the performance of LLMs on multimodal tasks and identify areas for improvement, which is essential for developing models that can effectively process and generate multimodal content.

To ensure the reliability and validity of benchmarking frameworks, it is essential to use high-quality datasets. \cite{ref201} is a benchmarking framework that provides a comprehensive evaluation of LLMs on code-related tasks, using a large-scale dataset of high-quality coding problems. This framework enables researchers to evaluate the performance of LLMs on code-related tasks and identify areas for improvement, which is vital for developing models that can be used in software development and other coding applications.

The development of benchmarking frameworks is an ongoing process, with new frameworks being proposed regularly. For example, \cite{ref202} is a benchmarking framework that evaluates the performance of LLMs on multimodal tasks, including vision-language tasks. This framework includes a comprehensive evaluation suite, which enables researchers to evaluate the performance of LLMs on various tasks and identify areas for improvement. The use of benchmarking frameworks like MMBench enables researchers to develop more accurate and efficient LLMs, which can be used in a wide range of applications, including natural language processing, computer vision, and multimodal tasks.

In conclusion, benchmarking frameworks are essential for evaluating the performance of large language models and identifying areas for improvement. The benchmarking frameworks discussed in this subsection provide a comprehensive evaluation of LLMs on various tasks and enable researchers to identify areas for improvement. The use of benchmarking frameworks has several advantages, including enabling researchers to evaluate the performance of LLMs on various tasks, identifying areas for improvement, and providing a fair and unbiased evaluation of LLMs. As the field of natural language processing continues to evolve, the use of benchmarking frameworks will become increasingly important for developing more accurate and efficient LLMs, which can be used in a wide range of applications. The insights gained from benchmarking frameworks will also inform the development of human-centric evaluation metrics, which are crucial for assessing the performance of LLMs in human-agent collaboration, as discussed in the following section.

\subsection{Human-Centric Evaluation Metrics}

Human-centric evaluation metrics are crucial for assessing the performance of agents in human-agent collaboration, as they focus on the user experience, social impact, and overall effectiveness of the interaction. These metrics go beyond traditional evaluation methods, which often prioritize technical performance over human-centered aspects \cite{ref203}. Human-centric evaluation metrics consider factors such as user satisfaction, trust, and perceived usefulness, providing a more comprehensive understanding of the agent's performance in real-world scenarios \cite{ref204}.

One key aspect of human-centric evaluation metrics is the assessment of user experience. This involves evaluating how users perceive the interaction with the agent, including factors such as usability, accessibility, and overall satisfaction \cite{ref205}. For instance, a study on human-AI collaboration in a social learning platform found that user-centric evaluation metrics, such as user satisfaction and perceived usefulness, were more effective in assessing the performance of the agent than traditional technical metrics \cite{ref206}. This highlights the importance of considering user experience in the evaluation of agents, as it can have a significant impact on the overall effectiveness of human-agent collaboration.

Another important aspect of human-centric evaluation metrics is the assessment of social impact. This involves evaluating the potential effects of the agent on users and society as a whole, including factors such as fairness, transparency, and accountability \cite{ref207}. For example, a study on the use of AI in hiring found that human-AI collaboration can lead to biased decision-making, highlighting the need for human-centric evaluation metrics to assess the social impact of such systems \cite{ref208}. This emphasizes the need for agents to be designed and developed with social impact in mind, in order to ensure that they are fair, transparent, and accountable.

Human-centric evaluation metrics can also be used to assess the effectiveness of agents in specific domains, such as education or healthcare. For instance, a study on the use of AI in education found that human-centric evaluation metrics, such as student engagement and motivation, were more effective in assessing the performance of the agent than traditional technical metrics \cite{ref209}. Similarly, a study on the use of AI in healthcare found that human-centric evaluation metrics, such as patient satisfaction and quality of life, were more effective in assessing the performance of the agent than traditional technical metrics \cite{ref210}. This demonstrates the versatility of human-centric evaluation metrics, which can be applied to a wide range of domains and applications.

In addition to these domain-specific applications, human-centric evaluation metrics can also be used to assess the overall effectiveness of agents in human-agent collaboration. This involves evaluating factors such as collaboration quality, communication effectiveness, and mutual understanding \cite{ref211}. For example, a study on human-AI collaboration in a game-playing scenario found that human-centric evaluation metrics, such as collaboration quality and mutual understanding, were more effective in assessing the performance of the agent than traditional technical metrics \cite{ref212}. This highlights the importance of considering the overall effectiveness of agents in human-agent collaboration, as it can have a significant impact on the success of the collaboration.

The use of human-centric evaluation metrics also raises important questions about the design and development of agents. For instance, how can agents be designed to prioritize human-centered aspects, such as user experience and social impact, over technical performance \cite{ref213}? How can human-centric evaluation metrics be integrated into the development process, and what are the implications for the design of agents and human-agent collaboration systems \cite{ref214}? These questions emphasize the need for a human-centered approach to the design and development of agents, which prioritizes user experience, social impact, and overall effectiveness.

To address these questions, researchers and developers can draw on a range of human-centric evaluation metrics and methodologies. For example, user studies and surveys can be used to assess user experience and satisfaction, while social impact assessments can be used to evaluate the potential effects of agents on users and society \cite{ref215}. Additionally, human-centric evaluation metrics can be integrated into the development process through the use of design thinking and human-centered design methodologies \cite{ref216}. This can help to ensure that agents are designed and developed with human-centered aspects in mind, leading to more effective and beneficial human-agent collaboration.

In conclusion, human-centric evaluation metrics are essential for assessing the performance of agents in human-agent collaboration. These metrics prioritize user experience, social impact, and overall effectiveness, providing a more comprehensive understanding of the agent's performance in real-world scenarios. By integrating human-centric evaluation metrics into the development process, researchers and developers can design and develop agents that prioritize human-centered aspects, leading to more effective and beneficial human-agent collaboration \cite{ref210}. As the field of human-AI collaboration continues to evolve, the use of human-centric evaluation metrics will play an increasingly important role in shaping the design and development of agents and human-agent collaboration systems \cite{ref217}.

Furthermore, the development of human-centric evaluation metrics is an ongoing process, with new metrics and methodologies being proposed and tested \cite{ref218}. For example, researchers have proposed the use of multimodal user behaviors, such as speech and gesture, to assess the human-likeness of conversational robots \cite{ref219}. Additionally, the use of large language models (LLMs) has been proposed as a means of evaluating the quality of meeting summaries, with the development of frameworks such as MESA \cite{ref220}. These developments highlight the ongoing need for innovative and effective human-centric evaluation metrics, which can be used to assess the performance of agents in human-agent collaboration.

The use of human-centric evaluation metrics also has important implications for the design of human-agent collaboration systems. For instance, the development of agents that prioritize human-centered aspects, such as user experience and social impact, will require the integration of human-centric evaluation metrics into the development process \cite{ref213}. This will involve the use of design thinking and human-centered design methodologies, as well as the development of new metrics and methodologies for assessing the performance of agents in human-agent collaboration \cite{ref216}. By prioritizing human-centered aspects, researchers and developers can create agents that are more effective, efficient, and beneficial for users.

In addition to these implications, the use of human-centric evaluation metrics also raises important questions about the future of human-AI collaboration. For example, how will the development of agents that prioritize human-centered aspects impact the nature of human-AI collaboration \cite{ref221}? How will the use of human-centric evaluation metrics shape the design and development of agents and human-agent collaboration systems \cite{ref217}? And what are the potential benefits and challenges of using human-centric evaluation metrics in the development of human-AI collaboration systems \cite{ref203}? These questions emphasize the need for ongoing research and development in the field of human-AI collaboration, as well as the importance of prioritizing human-centered aspects in the design and development of agents.

Overall, the use of human-centric evaluation metrics is essential for assessing the performance of agents in human-agent collaboration. These metrics prioritize user experience, social impact, and overall effectiveness, providing a more comprehensive understanding of the agent's performance in real-world scenarios. As the field of human-AI collaboration continues to evolve, the use of human-centric evaluation metrics will play an increasingly important role in shaping the design and development of agents and human-agent collaboration systems \cite{ref210}. By integrating human-centric evaluation metrics into the development process, researchers and developers can design and develop agents that prioritize human-centered aspects, leading to more effective and beneficial human-agent collaboration.

\section{Applications of Optimized Large Language Model-based Agents}

\subsection{Natural Language Processing Applications}

Natural language processing (NLP) has been a crucial area of research in artificial intelligence, and the emergence of large language models (LLMs) has significantly advanced the field \cite{ref222}. The applications of optimized large language model-based agents in NLP are diverse and continue to expand. In this subsection, we will delve into the applications of these agents in various NLP tasks, including text generation, language translation, and sentiment analysis, and explore their potential in advancing the field of NLP.

One of the primary applications of optimized large language model-based agents is text generation. These models can generate coherent and context-specific text based on a given prompt or input \cite{ref223}. For instance, \cite{ref223} demonstrated the use of fine-tuned transformer models for generating high-quality business names, highlighting the potential of these models in content creation and other applications. The study showed that larger models, which require more training time, yield better results for generating relatively short texts, such as business names.

Another significant application of optimized large language model-based agents is language translation. These models can learn to translate text from one language to another, given sufficient training data \cite{ref224}. For example, \cite{ref224} explored the use of BLOOMZ-3b, a large language model, for translating text from English to Indic languages like Hindi, Kannada, Malayalam, Tamil, and Telugu. The study demonstrated the effectiveness of using prompting and LoRA fine-tuning to improve the translation quality, highlighting the potential of these models in cross-lingual applications.

Sentiment analysis is a crucial NLP task that involves determining the sentiment or emotional tone of a given text \cite{ref222}. Optimized large language model-based agents have been applied to sentiment analysis tasks, achieving state-of-the-art results \cite{ref225}. For instance, \cite{ref225} evaluated the capabilities of large language models in performing various sentiment analysis tasks, from conventional sentiment classification to aspect-based sentiment analysis and multifaceted analysis of subjective texts. The study demonstrated that while large language models show satisfactory performance in simpler tasks, they lag behind in more complex tasks requiring deeper understanding or structured sentiment information.

In addition to these applications, optimized large language model-based agents have been used in other NLP tasks like question answering, text summarization, and named entity recognition \cite{ref226}. For example, \cite{ref226} compared the cross-lingual transfer capability of public small multilingual language models like XLM-R against English-centric large language models like Llama-3 in the context of sentiment analysis across English, Spanish, French, and Chinese. The study demonstrated that public small multilingual language models exhibit superior zero-shot cross-lingual performance relative to large language models. However, in few-shot cross-lingual settings, public large language models demonstrate enhanced adaptive potential.

The applications of optimized large language model-based agents in NLP are vast and continue to grow. These models have shown significant potential in advancing the field of NLP, and their applications are expected to expand to other areas like computer vision, robotics, and multimodal processing \cite{ref227}. The use of large language models in NLP has significant implications for applications like content creation, chatbots, language translation, and cross-lingual sentiment analysis. As the field of NLP continues to evolve, we can expect to see further advancements in the development and application of optimized large language model-based agents.

Furthermore, the development of large language models has also led to the creation of new benchmarks and evaluation metrics for NLP tasks \cite{ref31}. The development of new benchmarks and evaluation metrics is crucial for advancing the field of NLP and ensuring that large language models are evaluated fairly and accurately. Additionally, the development of large language models has also led to the creation of new techniques for improving the performance of these models \cite{ref228}. These techniques, such as parameter-efficient fine-tuning, pruning, and quantization, can significantly improve the performance of large language models, making them more efficient and effective in NLP tasks.

In conclusion, optimized large language model-based agents have been applied to various NLP tasks, including text generation, language translation, and sentiment analysis. These models have shown significant potential in advancing the field of NLP, and their applications are expected to expand to other areas like computer vision, robotics, and multimodal processing. As the field of NLP continues to evolve, we can expect to see further advancements in the development and application of optimized large language model-based agents, leading to new and innovative solutions to complex NLP tasks. The following subsection will explore the applications of optimized large language model-based agents in computer vision, highlighting their potential in advancing the field of computer vision and their implications for various applications.

\subsection{Computer Vision Applications}

The applications of optimized large language model-based agents in computer vision are vast and varied, with significant advancements in recent years. Building on the progress made in natural language processing, these agents have been applied to various computer vision tasks, including image recognition, object detection, and image generation. For instance, the Vision Transformer (ViT) has been shown to achieve state-of-the-art results in image classification tasks, demonstrating the potential of large language models in computer vision \cite{ref229}. Furthermore, the use of large language models in image recognition has also been explored in the context of few-shot learning, where the model is able to learn from a limited number of examples and generalize to new, unseen data \cite{ref230}.

In addition to image recognition, optimized large language model-based agents have also been applied to object detection tasks. Object detection is a critical task in computer vision, with numerous applications in areas such as autonomous vehicles, surveillance, and robotics \cite{ref231}. The use of large language models in object detection has been shown to improve the accuracy and efficiency of detection tasks, particularly in scenarios where the objects of interest are complex or have varying appearances \cite{ref231}. For example, the YOLO-RD model, which incorporates a Retriever-Dictionary module to retrieve features from a dictionary of visual models, has been shown to achieve state-of-the-art results in object detection tasks \cite{ref232}.

The applications of optimized large language model-based agents in computer vision also extend to image generation tasks. Image generation is a complex task that involves generating new images that are similar in style and content to a given set of images \cite{ref233}. The use of large language models in image generation has been shown to improve the quality and diversity of generated images, particularly in scenarios where the images of interest are complex or have varying styles \cite{ref233}. For instance, the MuLan model, which uses a multimodal large language model to generate images, has been shown to achieve state-of-the-art results in image generation tasks \cite{ref233}.

Moreover, optimized large language model-based agents have been applied to other computer vision tasks, such as image segmentation, image captioning, and visual question answering \cite{ref234}. The Uni-Perceiver v2 model, which is a generalist model that can be applied to a wide range of vision and vision-language tasks, has been shown to achieve state-of-the-art results in image segmentation and visual question answering tasks \cite{ref234}. Additionally, the use of large language models in multimodal learning has been explored, where the model learns from multiple sources of data, such as images, text, and audio \cite{ref235}.

The potential of optimized large language model-based agents in computer vision is significant, and it is likely that these agents will play a major role in the development of future computer vision systems \cite{ref20}. As the field of computer vision continues to evolve, it is likely that the applications of optimized large language model-based agents will continue to expand, leading to new and innovative solutions to complex computer vision tasks. The use of large language models in computer vision has the potential to improve the accuracy and efficiency of a wide range of computer vision tasks, and is likely to have a significant impact on various applications, including autonomous vehicles, surveillance, and robotics. The following section will explore the applications of optimized large language model-based agents in robotics and autonomous systems, highlighting their potential in advancing the field of robotics and autonomous systems.

\subsection{Robotics and Autonomous Systems}

The applications of optimized large language model-based agents in robotics and autonomous systems have shown tremendous potential in recent years, building on the advancements made in computer vision tasks. As discussed in the previous section, the use of large language models in computer vision has improved the accuracy and efficiency of tasks such as image recognition, object detection, and image generation. Similarly, in robotics and autonomous systems, these agents have been used in various tasks, including robotic manipulation, navigation, and human-robot interaction.

Robotic manipulation is a crucial aspect of robotics, where robots are required to manipulate objects in their environment. Optimized large language model-based agents have been used to improve robotic manipulation tasks, such as grasping and moving objects \cite{ref236}. These agents can learn from human feedback and adapt to new situations, making them more efficient and effective in robotic manipulation tasks. For example, \cite{ref237} presents a framework that combines the advantages of large language models with YOLO-based environmental perception to enable robots to autonomously make reasonable decisions and task planning based on given commands.

Navigation is another essential aspect of robotics, where robots need to navigate through their environment to reach a target location. Optimized large language model-based agents have been used to improve navigation tasks, such as mapping and localization \cite{ref238}. These agents can learn from human feedback and adapt to new situations, making them more efficient and effective in navigation tasks. For example, \cite{ref239} presents a system that uses large language models to interpret voice commands and generate sequential actions for tasks.

Human-robot interaction is a critical aspect of robotics, where robots need to interact with humans in a natural and intuitive way. Optimized large language model-based agents have been used to improve human-robot interaction tasks, such as dialogue systems and gesture recognition \cite{ref240}. These agents can learn from human feedback and adapt to new situations, making them more efficient and effective in human-robot interaction tasks. For example, \cite{ref241} presents a framework that exploits the inherent capabilities of pre-trained large language models, multimodal visual language models, and speech recognition models to decode high-level natural language conversations and semantic understanding of the robot's task environment.

In addition to these applications, optimized large language model-based agents have also been used in other areas of robotics and autonomous systems, such as robotic assembly and disassembly \cite{ref242}. These agents can learn from human feedback and adapt to new situations, making them more efficient and effective in robotic assembly and disassembly tasks.

The use of optimized large language model-based agents in robotics and autonomous systems has several advantages, including improved efficiency, effectiveness, and adaptability. These agents can learn from human feedback and adapt to new situations, making them more efficient and effective in various robotic tasks. Additionally, these agents can be used in a wide range of applications, from robotic manipulation and navigation to human-robot interaction and robotic assembly and disassembly.

However, there are also several challenges and limitations associated with the use of optimized large language model-based agents in robotics and autonomous systems. One of the main challenges is the need for large amounts of data and computational resources to train these agents. Additionally, these agents can be prone to errors and biases, which can affect their performance and reliability. Furthermore, the use of these agents in real-world applications requires careful consideration of safety and security concerns.

To address these challenges and limitations, researchers and developers are working on improving the efficiency, effectiveness, and adaptability of optimized large language model-based agents in robotics and autonomous systems. This includes developing new architectures and algorithms for these agents, as well as improving their ability to learn from human feedback and adapt to new situations. Additionally, researchers are working on developing new applications and use cases for these agents, such as robotic exploration and mapping \cite{ref243}.

In conclusion, the applications of optimized large language model-based agents in robotics and autonomous systems have shown tremendous potential in recent years, with a wide range of applications in robotic manipulation, navigation, human-robot interaction, and other areas. As these agents continue to evolve and improve, we can expect to see significant advances in the field of robotics and autonomous systems, enabling robots to perform complex tasks more efficiently and effectively. The future of optimized large language model-based agents in robotics and autonomous systems is promising, with potential applications in a wide range of areas, from robotic exploration and mapping to human-robot interaction and robotic assembly and disassembly. For example, \cite{ref244} presents a framework for decentralized multi-agent navigation, leveraging LLM-enabled communication and collaboration. Moreover, the use of optimized large language model-based agents in robotics and autonomous systems has the potential to revolutionize the way we interact with robots and autonomous systems, enabling more natural and intuitive human-robot interaction \cite{ref245}.

\section{Challenges and Limitations of Agent Optimization}

\subsection{Data Quality Challenges}

Data quality is a critical factor in the optimization of large language model-based agents, as it directly impacts the performance and reliability of these models. The challenges related to data quality, including data scarcity, noise, and bias, can significantly affect the optimization process and the overall effectiveness of large language model-based agents. In this subsection, we will delve into these challenges and their impact on the optimization of large language model-based agents, building on the discussion of scalability limitations in the previous section.

The amount of high-quality training data required to optimize large language models is substantial, and collecting such data can be time-consuming and expensive \cite{ref246}. Furthermore, the data collection process can be prone to errors, which can lead to noisy or biased data \cite{ref35}. Noisy or biased data can have a detrimental effect on the performance of large language model-based agents, leading to suboptimal results or even catastrophic failures. For instance, noise in the data can take many forms, including typos, grammatical errors, or inconsistencies in formatting \cite{ref53}, which can be introduced during the data collection process or can be inherent in the data itself.

Bias in the data is another significant challenge in the optimization of large language model-based agents. Bias can be introduced during the data collection process, or it can be inherent in the data itself \cite{ref247}. For example, if the training data is collected from a specific geographic region or demographic group, it may not be representative of the entire population, leading to biased results \cite{ref248}. Biased data can result in unfair or discriminatory outcomes, which can have severe consequences in real-world applications. The impact of data quality challenges on the optimization of large language model-based agents cannot be overstated, as poor data quality can lead to suboptimal results, reduced reliability, and even catastrophic failures \cite{ref40}.

To address these challenges, it is essential to use techniques such as data cleaning, data augmentation, and bias mitigation strategies \cite{ref249}. Data augmentation techniques, such as paraphrasing or text generation \cite{ref250}, can help increase the amount of training data available, reducing the impact of data scarcity. Additionally, data augmentation can also help reduce noise and bias in the data, by introducing new examples that are less prone to errors or biases \cite{ref251}. Another approach to addressing data quality challenges is to use bias mitigation strategies, such as debiasing or adversarial training \cite{ref252}, which can help reduce the impact of bias in the data and improve the fairness and transparency of large language model-based agents.

In conclusion, data quality challenges are a significant obstacle in the optimization of large language model-based agents. Addressing these challenges proactively, using techniques such as data cleaning, data augmentation, and bias mitigation strategies, is essential to achieving optimal results. By prioritizing data quality, we can develop more reliable, fair, and transparent large language model-based agents, that can be used in a wide range of applications, from natural language processing to decision-making and problem-solving \cite{ref31}. As we move forward to discuss the scalability limitations of optimizing large language model-based agents in the next section, it is crucial to consider the impact of data quality on the overall performance and reliability of these models.

\subsection{Scalability Limitations}

The scalability limitations of optimizing large language model-based agents are a significant concern, as these models require substantial computational resources and large amounts of data to achieve optimal performance \cite{ref34}. As discussed in the previous section, data quality challenges can significantly impact the optimization process, and scalability limitations can further exacerbate these challenges. The need for large amounts of computational resources and data is a major bottleneck in the development and deployment of these models, making it challenging to scale them up to meet the demands of real-world applications \cite{ref253}. Furthermore, the complexity of these models and the sheer size of the datasets required to train them make it difficult to optimize them for specific tasks or domains \cite{ref35}.

One of the primary scalability limitations of large language model-based agents is the need for significant computational resources \cite{ref173}. Training these models requires massive amounts of computational power, which can be costly and time-consuming \cite{ref17}. Moreover, the inference process also requires substantial computational resources, which can be a challenge for real-time applications \cite{ref254}. The computational requirements of these models can be mitigated to some extent by using distributed computing architectures \cite{ref255}, but this can add complexity to the system and require significant expertise to implement. Additionally, the use of techniques such as model pruning, quantization, and knowledge distillation can also help reduce the computational requirements of these models \cite{ref17}.

Another significant scalability limitation of large language model-based agents is the need for large amounts of data \cite{ref256}. These models require massive datasets to train and fine-tune, which can be challenging to obtain, especially for specific domains or tasks \cite{ref257}. The quality and diversity of the data are also crucial, as low-quality or biased data can negatively impact the performance of the model. Furthermore, the data requirements of these models can be mitigated to some extent by using data augmentation techniques \cite{ref33}, but this can add complexity to the system and require significant expertise to implement. The importance of high-quality data is closely related to the challenges discussed in the previous section, where data quality was identified as a critical factor in the optimization of large language model-based agents.

The scalability limitations of large language model-based agents also extend to the optimization process itself \cite{ref80}. Optimizing these models requires significant expertise and computational resources, which can be challenging to obtain, especially for small-scale applications \cite{ref258}. Moreover, the optimization process can be time-consuming and require significant trial and error, which can be a challenge for applications that require rapid development and deployment \cite{ref259}. The complexity of the optimization process is closely related to the challenges discussed in the following section, where the interpretability and explainability of large language model-based agents will be discussed.

In addition to the computational and data requirements, the scalability limitations of large language model-based agents also extend to the model architecture itself \cite{ref260}. The complexity of these models and the sheer size of the number of parameters required to achieve optimal performance can make it challenging to optimize them for specific tasks or domains \cite{ref261}. Furthermore, the model architecture can also impact the interpretability and explainability of the model, which can be a challenge for applications that require transparency and accountability \cite{ref262}. The importance of model architecture will be further discussed in the following section, where the interpretability and explainability of large language model-based agents will be explored in more detail.

In conclusion, the scalability limitations of optimizing large language model-based agents are a significant concern, requiring substantial computational resources, large amounts of data, and significant expertise to optimize and deploy \cite{ref11}. These limitations can be mitigated to some extent by using techniques such as distributed computing, data augmentation, and model pruning, but require significant expertise and computational resources to implement \cite{ref263}. Furthermore, the scalability limitations of these models can also impact the interpretability and explainability of the model, requiring significant expertise to optimize and deploy \cite{ref228}. As the field continues to evolve, it is essential to develop techniques and architectures that can mitigate these limitations and enable the widespread adoption of large language model-based agents in real-world applications \cite{ref264}.

\subsection{Interpretability and Explainability Challenges}

The interpretability and explainability of large language model-based agents are crucial aspects that need to be addressed to ensure the transparency and trustworthiness of these systems. As \cite{ref265} highlights, the lack of interpretability in large language models poses significant challenges, particularly in high-stakes applications where the consequences of incorrect decisions can be severe. This issue is closely related to the scalability limitations of large language model-based agents, as the complexity and size of these models can make it difficult to understand their internal mechanisms and decision-making processes.

One of the primary challenges in achieving interpretability is the complexity of large language models. As \cite{ref36} notes, the scale and complexity of these models make it difficult to understand their internal mechanisms. This complexity is further exacerbated by the fact that large language models are often trained on vast amounts of data, which can lead to a lack of transparency in their decision-making processes. To address this challenge, researchers have proposed various techniques, such as feature importance and model interpretability methods, to provide insights into the decision-making processes of large language models.

In addition to the complexity of large language models, the need for domain-specific knowledge is another significant challenge in achieving interpretability. As \cite{ref266} highlights, the interpretation of large language models requires domain-specific knowledge and expertise. This is particularly important in applications such as finance, where the consequences of incorrect decisions can be severe. To address this challenge, researchers have proposed the use of domain-specific lexicons and ontologies to provide a framework for interpreting large language models.

The use of attention mechanisms is another challenge in achieving interpretability. As \cite{ref267} notes, attention mechanisms can provide insights into the decision-making processes of large language models. However, the use of attention mechanisms can also lead to a lack of transparency, particularly if the attention weights are not properly calibrated. To address this challenge, researchers have proposed various techniques, such as attention visualization and attention regularization, to provide insights into the attention mechanisms of large language models.

Furthermore, the need for human-centered evaluation metrics is also an important challenge in achieving interpretability. As \cite{ref268} highlights, human-centered evaluation metrics are essential to evaluate the interpretability of large language models. This is particularly important in applications such as text-to-SQL semantic parsing, where the consequences of incorrect decisions can be severe. To address this challenge, researchers have proposed various human-centered evaluation metrics, such as user studies and surveys, to evaluate the interpretability of large language models.

The use of multimodal data is another challenge in achieving interpretability. As \cite{ref269} notes, multimodal data can provide rich insights into the decision-making processes of large language models. However, the use of multimodal data can also lead to a lack of transparency, particularly if the data is not properly calibrated. To address this challenge, researchers have proposed various techniques, such as multimodal attention and multimodal regularization, to provide insights into the decision-making processes of large language models.

In addition to these challenges, the need for explainability techniques is also an important aspect of achieving interpretability. As \cite{ref270} highlights, explainability techniques are essential to provide insights into the decision-making processes of large language models. This is particularly important in applications such as text classification, where the consequences of incorrect decisions can be severe. To address this challenge, researchers have proposed various explainability techniques, such as feature importance and model interpretability methods, to provide insights into the decision-making processes of large language models.

The use of reinforcement learning is another challenge in achieving interpretability. As \cite{ref271} notes, reinforcement learning can provide rich insights into the decision-making processes of large language models. However, the use of reinforcement learning can also lead to a lack of transparency, particularly if the rewards are not properly calibrated. To address this challenge, researchers have proposed various techniques, such as reward shaping and reward regularization, to provide insights into the decision-making processes of large language models.

Finally, the need for transparency in decision-making processes is also an important challenge in achieving interpretability. As \cite{ref272} highlights, transparency in decision-making processes is essential to foster trust and confidence in large language models. This is particularly important in applications such as decision-making, where the consequences of incorrect decisions can be severe. To address this challenge, researchers have proposed various techniques, such as model interpretability methods and feature importance, to provide insights into the decision-making processes of large language models.

In conclusion, the interpretability and explainability of large language model-based agents are crucial aspects that need to be addressed to ensure the transparency and trustworthiness of these systems. The challenges discussed above, including the complexity of large language models, the need for domain-specific knowledge, the use of attention mechanisms, the need for human-centered evaluation metrics, the use of multimodal data, the need for explainability techniques, the use of reinforcement learning, and the need for transparency in decision-making processes, highlight the importance of addressing interpretability and explainability in large language model-based agents. By addressing these challenges, researchers can develop more transparent and trustworthy large language model-based agents that can be used in a wide range of applications, including decision-making, text classification, and multimodal processing. As \cite{ref273} notes, the development of interpretable deep learning models is essential to bridge the gap between AI and human understanding, and to develop more transparent and trustworthy AI systems.

\section{Multimodal and Multitask Learning with Large Language Models}

\subsection{Multimodal Large Language Models}

Multimodal large language models (MLLMs) have gained significant attention in recent years due to their ability to process and understand multiple forms of data, such as text, images, and audio. These models are designed to learn representations that can capture the relationships between different modalities, enabling them to perform a wide range of tasks, including visual question answering, image captioning, and text-to-image synthesis. The architecture of MLLMs typically consists of vision and language encoders, which are used to extract features from images and text, respectively. The integration of these features is a key challenge in the development of MLLMs, and various approaches have been proposed to address this challenge, including the use of a single encoder to process multiple modalities, or the use of multiple encoders, each specialized to a particular modality \cite{ref274}.

The training of MLLMs typically involves the use of large datasets that contain multiple forms of data, such as text, images, and audio. These datasets are often constructed by combining existing datasets from different modalities, such as the combination of text datasets and image datasets to create a multimodal dataset \cite{ref275}. The models are then trained using a range of objectives, including masked language modeling, next sentence prediction, and image-text matching \cite{ref78}. One of the key benefits of MLLMs is their ability to perform a wide range of tasks, including visual question answering, image captioning, and text-to-image synthesis. For example, the MM1.5 model is able to perform visual question answering by using a vision encoder to extract features from images, and then combining these features with language features to generate answers to questions \cite{ref276}.

The use of MLLMs has also been explored in a range of applications, including robotics and computer vision. For example, the PaLM-E model uses a MLLM to enable a robot to understand and respond to natural language instructions, and to generate text descriptions of images \cite{ref21}. Similarly, the Valley model uses a MLLM to enable a robot to understand and respond to visual and linguistic inputs, and to generate text descriptions of images \cite{ref277}. In addition to their use in robotics and computer vision, MLLMs have also been explored in a range of other applications, including natural language processing and speech recognition. For example, the Discrete Multimodal Transformers model uses a MLLM to perform speech recognition and text-to-speech synthesis \cite{ref278}.

Despite the many benefits of MLLMs, there are also several challenges associated with their development and use. One of the key challenges is the need for large datasets that contain multiple forms of data, which can be difficult and expensive to construct \cite{ref279}. Another challenge is the need for specialized hardware and software to train and deploy MLLMs, which can be a barrier to entry for many researchers and practitioners \cite{ref280}. However, the benefits of MLLMs make them an exciting and important area of research, with the potential to enable a wide range of applications, including visual question answering, image captioning, and text-to-image synthesis \cite{ref281}.

The future of MLLMs is likely to involve the development of more advanced architectures and training methods, as well as the exploration of new applications and use cases. For example, the use of MLLMs in multimodal dialogue systems, which enable humans to interact with computers using multiple forms of input, such as text, images, and speech \cite{ref282}. Similarly, the use of MLLMs in multimodal sentiment analysis, which enables computers to understand and analyze human emotions and sentiments from multiple forms of input, such as text, images, and speech \cite{ref283}. As research in this area continues to evolve, we can expect to see even more innovative applications of MLLMs in the future, and the potential for these models to make a significant impact on the way we live and work \cite{ref284}.

In conclusion, MLLMs are a powerful tool for processing and understanding multiple forms of data, and have the potential to enable a wide range of applications, including visual question answering, image captioning, and text-to-image synthesis. The architecture and training methods of MLLMs are critical to their success, and involve the use of vision and language encoders, and the integration of multiple modalities. While there are several challenges associated with the development and use of MLLMs, the benefits of these models make them an exciting and important area of research \cite{ref285}. As we continue to develop and improve these models, we can expect to see even more impressive performance on a wide range of tasks, and their potential applications will continue to grow and expand into new areas.

\subsection{Applications of Multimodal Large Language Models}

Multimodal large language models have been increasingly applied to various tasks, showcasing their versatility and potential in handling multiple forms of input data. One of the primary applications of these models is in image-text retrieval, where the goal is to retrieve relevant images based on a given text query or vice versa \cite{ref286}. This task is crucial in many real-world applications, such as search engines, social media platforms, and e-commerce websites, where users often search for images using text queries. The success of multimodal large language models in image-text retrieval can be attributed to their ability to learn effective representations of both visual and textual data, allowing them to reason about the relationships between visual and textual information.

Another significant application of multimodal large language models is in visual question answering (VQA), which involves answering questions about an image \cite{ref287}. VQA is a challenging task that requires the model to understand both the visual content of the image and the natural language question. Multimodal large language models have shown impressive performance in VQA tasks, demonstrating their ability to reason about visual and textual information \cite{ref288}. Image captioning is another application where multimodal large language models have achieved remarkable results \cite{ref289}. Image captioning involves generating a natural language description of an image, which can be useful in various applications, such as image search, social media, and assistive technologies for visually impaired individuals.

In addition to these applications, multimodal large language models have also been applied to other tasks, such as visual grounding, referring expression comprehension, and multimodal sentiment analysis \cite{ref290}. Visual grounding involves identifying the objects or regions in an image that correspond to a given text description, while referring expression comprehension involves identifying the object or region in an image that corresponds to a given referring expression. Multimodal sentiment analysis involves analyzing the sentiment of text and images to determine the overall sentiment of a multimodal input. The applications of multimodal large language models are not limited to these tasks; they can also be used in more complex tasks, such as multimodal dialogue systems, visual storytelling, and multimodal machine translation \cite{ref291}.

The development of multimodal large language models has been driven by the availability of large-scale datasets that contain both visual and textual data \cite{ref292}. These datasets, such as Conceptual Captions, Visual Genome, and COCO, provide a rich source of training data for multimodal large language models. The use of these datasets has enabled researchers to develop models that can learn effective representations of both visual and textual data, leading to significant improvements in performance on a wide range of tasks. Furthermore, the development of multi-task models that can handle multiple tasks simultaneously has also been a key area of research \cite{ref293}. These models have been shown to achieve impressive performance on multiple tasks, demonstrating their ability to learn shared representations that are useful across multiple tasks.

As research in this area continues to evolve, we can expect to see even more innovative applications of multimodal large language models in the future. The potential for these models to improve the accessibility of visual content for people with visual impairments is also significant \cite{ref294}. For example, multimodal large language models can be used to generate audio descriptions of images, allowing people with visual impairments to access visual content. They can also be used to generate tactile graphics, which can be used to create tactile images that people with visual impairments can touch and explore. Overall, the applications of multimodal large language models are diverse and continue to grow as research in this area evolves.

In the future, we can expect to see even more impressive performance on a wide range of tasks, from image-text retrieval and visual question answering to multimodal dialogue systems and visual storytelling. The potential for multimodal large language models to revolutionize the way we interact with visual and textual data is significant, and their impact will be felt across a wide range of industries and applications \cite{ref295}. As these models continue to evolve, we can expect to see even more innovative applications in areas such as education, healthcare, and entertainment \cite{ref296}. Whether it's improving the accessibility of visual content, developing interactive learning systems, or generating interactive stories, multimodal large language models have the potential to make a significant impact on the way we live and work.

\subsection{Challenges and Limitations of Multimodal Learning}

Multimodal learning has emerged as a promising approach for integrating multiple sources of information to improve the performance of large language models. As discussed in the previous section, multimodal large language models have shown impressive performance in various tasks, including image-text retrieval, visual question answering, and image captioning \cite{ref297}. However, despite its potential, multimodal learning is not without its challenges and limitations. One of the primary challenges is the need for large-scale datasets that include multiple modalities, such as text, images, and audio \cite{ref297}. The collection and annotation of such datasets can be time-consuming and expensive, which can limit the widespread adoption of multimodal learning.

Another challenge is the complexity of integrating multiple modalities, which can require significant computational resources and expertise \cite{ref298}. The integration of multiple modalities can also lead to increased risk of hallucinations, which can negatively impact the performance of the model \cite{ref299}. Hallucinations can occur when the model generates text or other output that is not supported by the input data, which can be particularly problematic in applications where accuracy and reliability are critical. Furthermore, multimodal models can be prone to overfitting, particularly when the training data is limited or biased \cite{ref300}. Overfitting can cause the model to learn patterns and relationships that are specific to the training data, rather than generalizing to new, unseen data.

In addition to these challenges, multimodal learning can also be limited by the quality of the input data, which can be noisy, incomplete, or biased \cite{ref301}. Noisy or incomplete data can negatively impact the performance of the model, while biased data can perpetuate existing social and cultural biases. The use of multimodal data can also raise concerns about data privacy and security, particularly in applications where sensitive information is involved \cite{ref302}. Moreover, the evaluation of multimodal models can be complex, as it requires assessing the performance of the model across multiple modalities and tasks. The development of effective evaluation metrics and benchmarks is critical for advancing the field of multimodal learning and ensuring that models are fair, reliable, and accurate \cite{ref303}.

Despite these challenges and limitations, multimodal learning has the potential to revolutionize a wide range of applications, from natural language processing and computer vision to healthcare and education \cite{ref304}. The use of multimodal data can provide a more comprehensive and nuanced understanding of complex phenomena, which can lead to improved performance and accuracy in a wide range of tasks. To address the challenges and limitations of multimodal learning, researchers and practitioners can use a variety of techniques, such as data augmentation, transfer learning, and regularization \cite{ref305}. Additionally, techniques such as multimodal attention and fusion can be used to integrate multiple modalities and improve the performance of the model \cite{ref306}.

In the next section, we will discuss the applications of multimodal large language models in more detail, including their potential to improve the accessibility of visual content for people with visual impairments \cite{ref294}. We will also explore the future directions of multimodal learning and the potential for multimodal large language models to enable more sophisticated and human-like AI systems \cite{ref307}. By advancing the field of multimodal learning, we can develop more effective and reliable models that can interact with humans in a more natural and intuitive way, leading to significant improvements in a wide range of applications.

\section{Ethics and Safety Considerations for Large Language Models}

\subsection{Bias and Fairness in Large Language Models}

Bias and fairness in large language models are critical issues that have gained significant attention in recent years, particularly in the context of their potential to generate harmful or toxic content. The emergence of large language models (LLMs) \cite{ref32} has transformed the natural language processing (NLP) landscape, enabling a wide range of applications, from language translation to text generation. However, these models can perpetuate and amplify biases present in their training data, leading to unfair outcomes and potential harm to marginalized groups \cite{ref308}. This is closely related to the safety concerns discussed earlier, as biased models can generate content that is not only harmful but also discriminatory.

The sources of bias in LLMs are multifaceted and complex, and can be attributed to various factors, including the training data, model architecture, and training objectives. One primary source is the training data itself, which can reflect societal biases and stereotypes \cite{ref309}. For instance, if a model is trained on a dataset that contains more texts written by men than women, it may learn to associate certain topics or styles with men, leading to biased representations \cite{ref310}. Another source of bias is the model's architecture and training objectives, which can inadvertently introduce biases during the learning process \cite{ref311}. These biases can have severe consequences, including the perpetuation of stereotypes, reinforcement of existing power dynamics, and even discriminatory outcomes \cite{ref312}.

The impact of bias on marginalized groups can be severe and far-reaching, and is closely tied to the safety concerns related to LLMs. Biased LLMs can perpetuate stereotypes, reinforce existing power dynamics, and even lead to discriminatory outcomes \cite{ref312}. For example, a biased language model may be more likely to generate texts that reinforce negative stereotypes about certain racial or ethnic groups, contributing to a toxic online environment \cite{ref313}. Moreover, biased models can also lead to unfair treatment of individuals based on their demographic characteristics, such as gender, age, or disability \cite{ref314}. To mitigate these risks, it is essential to develop methods for detecting and mitigating bias in LLMs, which can be used in conjunction with methods for detecting and mitigating harmful content.

Techniques for mitigating bias in LLMs are an active area of research, and can be used to improve the safety and fairness of these models. One approach is to use debiasing methods, such as data preprocessing, regularization techniques, or adversarial training \cite{ref315}. Another approach is to use fairness-aware optimization methods, which can help to reduce bias by optimizing the model's parameters to minimize fairness metrics \cite{ref316}. Additionally, researchers have proposed using diverse and representative training data, as well as techniques such as data augmentation and transfer learning, to reduce bias \cite{ref317}. These techniques can be used to improve the safety and fairness of LLMs, and can be used in conjunction with methods for detecting and mitigating harmful content.

Despite these efforts, mitigating bias in LLMs remains a challenging task, particularly in the context of their potential to generate harmful or toxic content. One reason is that bias can be subtle and context-dependent, making it difficult to detect and measure \cite{ref318}. Another reason is that bias can be deeply ingrained in the model's representations, requiring significant changes to the model's architecture or training objectives \cite{ref252}. To address these challenges, researchers have proposed using a range of evaluation metrics and benchmarks to assess bias in LLMs \cite{ref319}. These metrics can help to identify biased models and track progress over time, and can be used to improve the safety and fairness of LLMs.

In conclusion, bias and fairness in large language models are critical issues that require careful attention and mitigation, particularly in the context of their potential to generate harmful or toxic content. The sources of bias are complex and multifaceted, and the impact of bias on marginalized groups can be severe. Techniques for mitigating bias are an active area of research, and can be used to improve the safety and fairness of LLMs. However, mitigating bias in LLMs remains a challenging task, and further research is needed to develop effective and scalable solutions \cite{ref320}. Ultimately, addressing bias and fairness in LLMs requires a multidisciplinary approach that involves researchers, practitioners, and stakeholders from diverse backgrounds and perspectives \cite{ref321}. By developing more robust and transparent LLMs, as well as more effective methods for detecting and mitigating bias, we can work towards creating more trustworthy and reliable models that can be used to benefit society.

\subsection{Safety and Harmful Content}

The safety concerns related to large language models (LLMs) are a pressing issue, as these models have the potential to generate harmful or toxic content \cite{ref322}. This can be particularly problematic when LLMs are used in applications such as text generation, chatbots, or language translation, where the generated content can be seen by a large audience \cite{ref323}. The generation of harmful content can have serious consequences, including the spread of misinformation, the promotion of hate speech or violence, and the perpetuation of harmful stereotypes \cite{ref324}. As discussed in the previous section, bias and fairness in LLMs are closely related to safety concerns, and mitigating these risks is essential to ensure the trustworthy and reliable deployment of LLMs.

To mitigate these risks, it is essential to develop methods for detecting and mitigating harmful content generated by LLMs \cite{ref325}. One approach is to use automated content moderation tools, which can detect and remove harmful content from online platforms \cite{ref326}. However, these tools are not foolproof and can be evaded by sophisticated attackers \cite{ref327}. Another approach is to use human evaluators to assess the safety of LLM-generated content \cite{ref328}. Human evaluators can provide more nuanced and context-specific assessments of content safety, but this approach can be time-consuming and expensive \cite{ref329}.

To address these challenges, researchers have proposed various methods for improving the safety of LLMs, including fine-tuning, pruning, and knowledge distillation \cite{ref330}. Fine-tuning involves adjusting the model's parameters to better align with safety goals, while pruning involves removing or modifying parts of the model that are prone to generating harmful content \cite{ref331}. Knowledge distillation involves training a smaller model to mimic the behavior of a larger, safer model \cite{ref332}. Additionally, researchers have proposed using multimodal approaches, which combine text and image inputs to improve the safety and accuracy of LLMs \cite{ref333}.

The development of more robust and transparent LLMs is also crucial for improving safety \cite{ref334}. This can help to identify and mitigate potential safety risks, as well as provide more trustworthy and reliable outputs \cite{ref335}. Furthermore, it is essential to consider the broader social and cultural context in which LLMs are deployed \cite{ref336}. This includes addressing issues related to bias, fairness, and inclusivity, as well as ensuring that LLMs are aligned with human values and safety goals \cite{ref337}. By taking a more holistic and nuanced approach to LLM safety, we can work towards developing more trustworthy and reliable models that can be used to benefit society \cite{ref338}.

In conclusion, the safety concerns related to large language models are a pressing issue that requires ongoing attention and research \cite{ref339}. By developing more robust and transparent LLMs, as well as more effective content moderation tools and strategies, we can work towards mitigating the risks associated with harmful content generation \cite{ref340}. Additionally, it is essential to consider the broader social and cultural context in which LLMs are deployed, and to ensure that these models are aligned with human values and safety goals \cite{ref341}. The development of transparent and explainable LLMs, as discussed in the following section, will also be crucial in addressing the safety concerns and risks associated with these models, and will have significant implications for the development of more advanced AI systems \cite{ref342}.

\subsection{Transparency and Explainability}

Transparency and explainability are essential components in the development and deployment of large language models (LLMs), particularly in light of the safety concerns and risks associated with these models, as discussed in the previous section \cite{ref265}. As LLMs become increasingly integrated into various aspects of our lives, understanding how they make decisions and arrive at their outputs is crucial for building trust, ensuring accountability, and identifying potential biases or errors \cite{ref343}. The importance of transparency and explainability in LLMs can be seen in their ability to provide insights into the decision-making processes of these complex models, which is vital for high-stakes applications such as healthcare, finance, and education \cite{ref36}.

One of the primary techniques for achieving transparency and explainability in LLMs is through the use of model interpretability methods, which can help identify the most important input features that contribute to the model's output \cite{ref270}. For instance, techniques such as feature attribution and model explainability can provide valuable insights into the decision-making process, allowing developers and users to understand how the model arrives at its decisions \cite{ref265}. Another approach to enhancing transparency and explainability in LLMs is through the use of natural language explanations, which involve generating human-readable text that describes the reasoning process behind the model's output \cite{ref343}.

In addition to these techniques, researchers have also explored the use of multimodal explanations, which involve generating explanations that incorporate multiple forms of media, such as text, images, and audio \cite{ref269}. Multimodal explanations have been shown to be effective in providing a more comprehensive understanding of the decision-making process of LLMs, particularly in tasks that involve multiple forms of input data \cite{ref344}. Furthermore, the development of transparent and explainable LLMs has significant implications for a wide range of applications, including healthcare, finance, and education, where the ability to understand and trust the decision-making process of these models is critical \cite{ref345}.

Despite the progress made in developing techniques for transparency and explainability in LLMs, there are still several challenges that need to be addressed, including the complexity of LLMs, which can make it difficult to develop interpretability methods that can effectively capture the decision-making process of these models \cite{ref346}. Another challenge is the need for more robust evaluation metrics, which can effectively assess the quality and accuracy of explanations generated by LLMs \cite{ref347}. To address these challenges, researchers have proposed several approaches, including the development of more advanced interpretability methods, such as attention mechanisms and layer-wise relevance propagation \cite{ref348}.

In conclusion, transparency and explainability are critical components in the development and deployment of large language models, and are essential for building trust, ensuring accountability, and identifying potential biases or errors \cite{ref349}. By developing more transparent and explainable LLMs, researchers can create models that are not only more accurate and reliable but also more trustworthy and accountable, which is critical for the safe and effective deployment of these models in a wide range of applications \cite{ref350}. The development of transparent and explainable LLMs will be crucial in addressing the safety concerns and risks associated with these models, and will have significant implications for the development of more advanced AI systems, such as autonomous vehicles and robots \cite{ref342}.

\section{Conclusion and Future Directions}

\subsection{Summary of Key Findings}

This subsection provides a comprehensive summary of the key findings from the survey, encompassing the current state of research in large language models, optimization techniques, and applications. The emergence of large language models (LLMs) \cite{ref33} has revolutionized the field of natural language processing, enabling models to achieve state-of-the-art performance in various tasks. The development of LLMs has been fueled by advances in deep learning technology, which provides new opportunities for the construction and application of these models \cite{ref33}. As a result, LLMs have become a crucial component in many AI systems, with applications ranging from natural language processing to computer vision and robotics.

One of the primary challenges associated with LLMs is their substantial computational and memory requirements, hindering their deployment in resource-constrained scenarios \cite{ref351}. To address this issue, researchers have explored various optimization techniques, including model compression, pruning, quantization, and knowledge distillation \cite{ref12}. These techniques aim to reduce the computational and memory demands of LLMs while maintaining their performance. For instance, \cite{ref86} presents a learning-based framework for mask selection, which enables the efficient pruning of LLMs. Furthermore, the development of efficient inference methods for LLMs \cite{ref11} has also been a key area of research, enabling the deployment of these models in real-time applications.

The applications of LLMs are diverse and widespread, ranging from natural language processing tasks such as text generation, language translation, and sentiment analysis \cite{ref44} to computer vision tasks like image recognition and object detection \cite{ref25}. LLMs have also been applied in robotics and autonomous systems, enabling robots to understand and generate human-like language. Additionally, \cite{ref352} highlights the potential of LLMs in code-related applications, such as code generation and completion. The survey also emphasizes the importance of evaluating the performance of LLMs using various metrics and benchmarks, such as \cite{ref201}, which provides a comprehensive evaluation framework for LLMs in code-related tasks.

Moreover, the survey discusses the challenges and limitations associated with LLMs, including issues related to data quality, scalability, interpretability, and robustness. These challenges must be addressed to ensure the reliable and efficient deployment of LLMs in real-world applications. In terms of future research directions, the survey suggests that there is a need for more efficient training methods and better evaluation metrics for LLMs. Additionally, \cite{ref4} emphasizes the importance of expanding the knowledge of LLMs, enabling them to adapt to new and diverse knowledge domains. \cite{ref353} highlights the need for sustainable NLP practices, emphasizing the importance of reducing the environmental impact of LLMs.

Overall, the survey provides a comprehensive overview of the current state of research in LLMs, highlighting their potential applications, optimization techniques, and challenges. The survey emphasizes the importance of efficient inference methods, evaluation metrics, and knowledge expansion techniques for LLMs, and suggests that there is a need for more research into the challenges and limitations associated with these models. By addressing these challenges and limitations, researchers can ensure the reliable and efficient deployment of LLMs in real-world applications, enabling them to achieve their full potential. As the field of LLMs continues to evolve, it is essential to explore new research directions, such as the application of LLMs in multimodal and multitask learning scenarios \cite{ref8}, and the development of more efficient algorithmic techniques for LLMs \cite{ref354}.

The survey also highlights the importance of exploring new applications and domains for LLMs, such as scientific research and code-related tasks. For instance, \cite{ref355} demonstrates the potential of LLMs in scientific research, enabling them to generate hypotheses and design experiments. Additionally, \cite{ref1} provides a detailed overview of the current state of research in LLMs, highlighting their architecture, training methods, and applications. \cite{ref356} emphasizes the potential of LLMs as embedding models, enabling them to be used in a wide range of NLP tasks. By continuing to advance the field of LLMs, researchers can unlock new possibilities for AI systems and enable them to achieve their full potential.

In conclusion, the survey provides a comprehensive summary of the key findings in the field of LLMs, highlighting their potential applications, optimization techniques, and challenges. The survey emphasizes the importance of efficient inference methods, evaluation metrics, and knowledge expansion techniques for LLMs, and suggests that there is a need for more research into the challenges and limitations associated with these models. As the field of LLMs continues to evolve, it is essential to address these challenges and limitations, enabling the reliable and efficient deployment of LLMs in real-world applications, and unlocking new possibilities for AI systems.

\subsection{Future Research Directions}

Future research directions in the optimization of large language model-based agents are vast and diverse, with several potential areas for improvement and emerging trends. Building on the existing literature, including \cite{ref357}, \cite{ref358}, and \cite{ref228}, several key areas of focus can be identified. One of the primary areas of focus is the development of more efficient and effective optimization techniques, such as fine-tuning and pruning methods \cite{ref104}, which can help reduce the computational costs and improve the performance of large language models, making them more suitable for real-world applications.

Another area of research is the exploration of new architectures and models, such as multimodal large language models and multi-agent systems \cite{ref359}, which have the potential to improve the performance and capabilities of large language models, enabling them to handle more complex tasks and scenarios. For instance, multimodal large language models can process and generate multiple forms of data, such as text, images, and audio, while multi-agent systems can facilitate collaboration and communication between multiple agents. The use of reinforcement learning and other machine learning techniques is also an active area of research \cite{ref165}, which can help improve the performance of large language models by enabling them to learn from feedback and adapt to new situations.

In addition to these areas, the development of more advanced evaluation metrics and benchmarks is necessary to assess the performance of large language models and identify areas for improvement \cite{ref351}. The application of large language models in various domains, such as natural language processing, computer vision, and robotics, is also an exciting area of research \cite{ref20}, which has the potential to revolutionize various industries and applications, such as language translation, text summarization, and chatbots. Furthermore, the development of more transparent and explainable large language models is essential to build trust and understanding in these models \cite{ref360}.

The integration of large language models with other AI technologies, such as computer vision and robotics, is also an emerging trend \cite{ref49}, which can enable the development of more sophisticated and capable AI systems that can perceive, understand, and interact with their environment. Moreover, the use of large language models in multimodal settings, such as human-computer interaction and human-robot interaction, is an area of research that can lead to more natural and intuitive interfaces \cite{ref25}. The development of more efficient and scalable large language models is also a critical area of research \cite{ref361}, which can be achieved through the use of techniques such as model pruning, knowledge distillation, and quantization.

As the field of large language models continues to evolve, it is essential to explore new applications and domains for these models, such as scientific research and code-related tasks \cite{ref355}. The development of more advanced optimization techniques, such as gradient-based optimization and evolutionary algorithms, is also an area of research that can help improve the performance of large language models \cite{ref362}. These techniques can help adapt large language models to new situations and improve their performance in various tasks and applications. Moreover, the use of large language models in multi-agent systems and collaborative learning scenarios is an emerging trend that can lead to more sophisticated and capable AI systems \cite{ref363}.

In conclusion, the optimization of large language model-based agents is a rapidly evolving field with numerous potential areas for improvement and emerging trends. The development of more efficient and effective optimization techniques, new architectures and models, and the application of large language models in various domains and scenarios are all exciting areas of research that can lead to more practical and effective AI systems. Additionally, the exploration of new applications and domains, the development of more advanced evaluation metrics and benchmarks, and the integration of large language models with other AI technologies are all critical areas of research that can help advance the field of AI and lead to more sophisticated and capable AI systems \cite{ref364}.

\subsection{Recommendations for Practitioners and Researchers}

As researchers and practitioners continue to push the boundaries of large language model-based agents, it is essential to provide recommendations for optimizing these models. Based on the existing literature, including \cite{ref357}, \cite{ref358}, and \cite{ref228}, several best practices and potential pitfalls can be identified.

To begin with, fine-tuning large language models requires careful consideration of the size of the training dataset and the complexity of the task at hand, as shown in \cite{ref80}. While fine-tuning with as few as 200 samples can improve model accuracy, the law of diminishing returns applies, and additional data beyond a certain point yields minimal gains. Therefore, practitioners should carefully evaluate the trade-off between dataset size and model performance.

The choice of optimization technique is also crucial, as techniques such as \cite{ref165} and \cite{ref365} have been shown to outperform traditional methods. However, these techniques may require significant computational resources and expertise to implement. As such, researchers and practitioners should carefully consider the strengths and weaknesses of each technique and select the most suitable approach for their specific use case.

In addition to selecting the right optimization technique, the importance of evaluation metrics and benchmarks cannot be overstated. As demonstrated in \cite{ref173}, the use of inappropriate evaluation metrics can lead to suboptimal performance and inefficient use of resources. Therefore, practitioners should carefully select evaluation metrics that align with their specific goals and objectives, and consider using multiple metrics to provide a comprehensive understanding of model performance.

Moreover, the potential for large language models to be used in a variety of applications, including \cite{ref49} and \cite{ref366}, highlights the need for researchers and practitioners to consider the broader implications of their work. This includes ensuring that models are fair, transparent, and accountable, and that they do not perpetuate existing biases or inequalities.

To avoid common pitfalls, researchers and practitioners should be aware of the risk of overfitting or underfitting, as shown in \cite{ref86}. To mitigate this risk, practitioners should consider using techniques such as regularization, dropout, or early stopping, and carefully evaluate model performance on a hold-out test set. Additionally, the failure to consider the limitations and biases of large language models can lead to vulnerabilities and attacks, as demonstrated in \cite{ref367}.

In conclusion, the optimization of large language model-based agents is a complex and multifaceted challenge that requires careful consideration of a range of factors. By following best practices, avoiding common pitfalls, and staying up-to-date with the latest research and developments, researchers and practitioners can create more effective, efficient, and robust agents that can be used to benefit society as a whole. As the field continues to evolve, it is essential to prioritize the development of fair, transparent, and accountable models that can be used to benefit society as a whole, and to consider the broader implications of large language models in a variety of applications.

Ultimately, the optimization of large language model-based agents has the potential to revolutionize a wide range of applications, from natural language processing to computer vision and robotics. As shown in \cite{ref368}, the use of large language models can lead to significant advancements in various fields. However, to unlock the full potential of these models, it is essential to prioritize the development of fair, transparent, and accountable models that can be used to benefit society as a whole. Future research directions may include the development of more efficient and effective optimization techniques, such as \cite{ref369} and \cite{ref370}, and the use of large language models in a variety of applications, including \cite{ref371} and \cite{ref372}.


\begin{thebibliography}{372}
\addcontentsline{toc}{section}{References}
\bibitem{ref1} A Comprehensive Overview of Large Language Models. arXiv:2307.06435, 2023. \url{https://arxiv.org/abs/2307.06435}
\bibitem{ref2} Efficient Large Language Models A Survey. arXiv:2312.03863, 2023. \url{https://arxiv.org/abs/2312.03863}
\bibitem{ref3} Recent Advances in Large Language Models for Healthcare. doi:10.3390/biomedinformatics4020062, 2024. \url{https://doi.org/10.3390/biomedinformatics4020062}
\bibitem{ref4} Bring Your Own Knowledge: A Survey of Methods for LLM Knowledge Expansion. arXiv:2502.12598, 2025. \url{https://arxiv.org/abs/2502.12598}
\bibitem{ref5} Large Language Models for Code Analysis: Do LLMs Really Do Their Job?. arXiv:2310.12357, 2023. \url{https://arxiv.org/abs/2310.12357}
\bibitem{ref6} Easy Problems That LLMs Get Wrong. arXiv:2405.19616, 2024. \url{https://arxiv.org/abs/2405.19616}
\bibitem{ref7} Through the Prism of Culture: Evaluating LLMs' Understanding of Indian Subcultures and Traditions. arXiv:2501.16748, 2025. \url{https://arxiv.org/abs/2501.16748}
\bibitem{ref8} Trends in Integration of Knowledge and Large Language Models: A Survey and Taxonomy of Methods, Benchmarks, and Applications. arXiv:2311.05876, 2023. \url{https://arxiv.org/abs/2311.05876}
\bibitem{ref9} A Survey on Mixture of Experts in Large Language Models. doi:10.1109/tkde.2025.3554028, 2025. \url{https://doi.org/10.1109/tkde.2025.3554028}
\bibitem{ref10} ChatGPT in society: emerging issues. doi:10.3389/frai.2023.1130913, 2023. \url{https://doi.org/10.3389/frai.2023.1130913}
\bibitem{ref11} A Survey on Efficient Inference for Large Language Models. arXiv:2404.14294, 2024. \url{https://arxiv.org/abs/2404.14294}
\bibitem{ref12} A Survey on Model Compression for Large Language Models. arXiv:2308.07633, 2023. \url{https://arxiv.org/abs/2308.07633}
\bibitem{ref13} Exploring Multilingual Probing in Large Language Models: A Cross-Language Analysis. arXiv:2409.14459, 2024. \url{https://arxiv.org/abs/2409.14459}
\bibitem{ref14} Security and Privacy Challenges of Large Language Models: A Survey. doi:10.1145/3712001, 2025. \url{https://doi.org/10.1145/3712001}
\bibitem{ref15} Surveying (Dis)Parities and Concerns of Compute Hungry NLP Research. arXiv:2306.16900, 2023. \url{https://arxiv.org/abs/2306.16900}
\bibitem{ref16} LARGE LANGUAGE MODELS: BUSINESS APPLICATIONS AND DEVELOPMENT PROSPECTS. doi:10.26906/sunz.2025.1.129-133, 2025. \url{https://doi.org/10.26906/sunz.2025.1.129-133}
\bibitem{ref17} Make Large Language Models Efficient: A Review. doi:10.1109/access.2025.3605110, 2025. \url{https://doi.org/10.1109/access.2025.3605110}
\bibitem{ref18} Several categories of Large Language Models (LLMs): A Short Survey. doi:10.22214/ijraset.2023.54677, 2023. \url{https://doi.org/10.22214/ijraset.2023.54677}
\bibitem{ref19} Two Heads Are Better Than One: Collaborative LLM Embodied Agents for Human-Robot Interaction. arXiv:2411.16723, 2024. \url{https://arxiv.org/abs/2411.16723}
\bibitem{ref20} Large Language Models Meet Computer Vision: A Brief Survey. arXiv:2311.16673, 2023. \url{https://arxiv.org/abs/2311.16673}
\bibitem{ref21} PaLM-E: An Embodied Multimodal Language Model. arXiv:2303.03378, 2023. \url{https://arxiv.org/abs/2303.03378}
\bibitem{ref22} Inner Monologue: Embodied Reasoning through Planning with Language Models. arXiv:2207.05608, 2022. \url{https://arxiv.org/abs/2207.05608}
\bibitem{ref23} Do As I Can, Not As I Say: Grounding Language in Robotic Affordances. arXiv:2204.01691, 2022. \url{https://arxiv.org/abs/2204.01691}
\bibitem{ref24} LLMs for Coding and Robotics Education. arXiv:2402.06116, 2024. \url{https://arxiv.org/abs/2402.06116}
\bibitem{ref25} Programming Language Models in Multilingual Settings. doi:10.1145/3639478.3639787, 2024. \url{https://doi.org/10.1145/3639478.3639787}
\bibitem{ref26} Scientific Large Language Models: A Survey on Biological \& Chemical Domains. doi:10.1145/3715318, 2025. \url{https://doi.org/10.1145/3715318}
\bibitem{ref27} Large language models and their applications in bioinformatics. doi:10.1016/j.csbj.2024.09.031, 2024. \url{https://doi.org/10.1016/j.csbj.2024.09.031}
\bibitem{ref28} Developing Interactive Tourism Planning: A Dialogue Robot System Powered by a Large Language Model. arXiv:2312.13545, 2023. \url{https://arxiv.org/abs/2312.13545}
\bibitem{ref29} Alchemist: LLM-Aided End-User Development of Robot Applications. doi:10.1145/3610977.3634969, 2024. \url{https://doi.org/10.1145/3610977.3634969}
\bibitem{ref30} Training a T5 Using Lab-sized Resources. arXiv:2208.12097, 2022. \url{https://arxiv.org/abs/2208.12097}
\bibitem{ref31} Evaluations of Large Language Models a Bibliometric analysis. doi:10.33022/ijcs.v13i1.3767, 2024. \url{https://doi.org/10.33022/ijcs.v13i1.3767}
\bibitem{ref32} A Survey on Fairness in Large Language Models. arXiv:2308.10149, 2023. \url{https://arxiv.org/abs/2308.10149}
\bibitem{ref33} Research on the Application and Optimization Strategies of Deep Learning in Large Language Models. doi:10.53469/jtpes.2024.04(05).12, 2024. \url{https://doi.org/10.53469/jtpes.2024.04(05).12}
\bibitem{ref34} Scaling Laws for BERT in Low-Resource Settings. doi:10.18653/v1/2023.findings-acl.492, 2023. \url{https://doi.org/10.18653/v1/2023.findings-acl.492}
\bibitem{ref35} To Repeat or Not To Repeat: Insights from Scaling LLM under Token-Crisis. arXiv:2305.13230, 2023. \url{https://arxiv.org/abs/2305.13230}
\bibitem{ref36} From Understanding to Utilization: A Survey on Explainability for Large Language Models. arXiv:2401.12874, 2024. \url{https://arxiv.org/abs/2401.12874}
\bibitem{ref37} Shortcut Learning Explanations for Deep Natural Language Processing: A Survey on Dataset Biases. doi:10.1109/access.2024.3360306, 2024. \url{https://doi.org/10.1109/access.2024.3360306}
\bibitem{ref38} NormXLogit: The Head-on-Top Never Lies. arXiv:2411.16252, 2024. \url{https://arxiv.org/abs/2411.16252}
\bibitem{ref39} On the Impossible Safety of Large AI Models. arXiv:2209.15259, 2022. \url{https://arxiv.org/abs/2209.15259}
\bibitem{ref40} Assessing LLMs for High Stakes Applications. doi:10.1145/3639477.3639720, 2024. \url{https://doi.org/10.1145/3639477.3639720}
\bibitem{ref41} Multilingual European Language Models: Benchmarking Approaches and Challenges. arXiv:2502.12895, 2025. \url{https://arxiv.org/abs/2502.12895}
\bibitem{ref42} Evaluating Large Language Models for Generalization and Robustness via Data Compression. arXiv:2402.00861, 2024. \url{https://arxiv.org/abs/2402.00861}
\bibitem{ref43} Curriculum: A Broad-Coverage Benchmark for Linguistic Phenomena in Natural Language Understanding. arXiv:2204.06283, 2022. \url{https://arxiv.org/abs/2204.06283}
\bibitem{ref44} Large Language Models: A Survey. arXiv:2402.06196, 2024. \url{https://arxiv.org/abs/2402.06196}
\bibitem{ref45} Language Model Fine-Tuning on Scaled Survey Data for Predicting Distributions of Public Opinions. arXiv:2502.16761, 2025. \url{https://arxiv.org/abs/2502.16761}
\bibitem{ref46} TPTU Large Language Model-based AI Agents for Task Planning and Tool Usage. arXiv:2308.03427, 2023. \url{https://arxiv.org/abs/2308.03427}
\bibitem{ref47} Sub-goal Distillation: A Method to Improve Small Language Agents. arXiv:2405.02749, 2024. \url{https://arxiv.org/abs/2405.02749}
\bibitem{ref48} A Survey on the Memory Mechanism of Large Language Model-based Agents. doi:10.1145/3748302, 2024. \url{https://doi.org/10.1145/3748302}
\bibitem{ref49} Web Agents with World Models: Learning and Leveraging Environment Dynamics in Web Navigation. arXiv:2410.13232, 2024. \url{https://arxiv.org/abs/2410.13232}
\bibitem{ref50} Large Language Model-Based Agents for Software Engineering: A Survey. arXiv:2409.02977, 2024. \url{https://arxiv.org/abs/2409.02977}
\bibitem{ref51} A Survey of LLM-based Agents in Medicine: How far are we from Baymax?. arXiv:2502.11211, 2025. \url{https://arxiv.org/abs/2502.11211}
\bibitem{ref52} Practical Applications of Large Language Models in Enterprise-Level Applications. doi:10.54097/4hnrtz02, 2025. \url{https://doi.org/10.54097/4hnrtz02}
\bibitem{ref53} EDGE: Efficient Data Selection for LLM Agents via Guideline Effectiveness. arXiv:2502.12494, 2025. \url{https://arxiv.org/abs/2502.12494}
\bibitem{ref54} The Emerged Security and Privacy of LLM Agent: A Survey with Case Studies. arXiv:2407.19354, 2024. \url{https://arxiv.org/abs/2407.19354}
\bibitem{ref55} Attention Is All You Need. arXiv:1706.03762, 2017. \url{https://arxiv.org/abs/1706.03762}
\bibitem{ref56} RecurFormer: Not All Transformer Heads Need Self-Attention. arXiv:2410.12850, 2024. \url{https://arxiv.org/abs/2410.12850}
\bibitem{ref57} Exploring the Limits of Language Modeling. arXiv:1602.02410, 2016. \url{https://arxiv.org/abs/1602.02410}
\bibitem{ref58} A Survey on Dynamic Neural Networks for Natural Language Processing. doi:10.18653/v1/2023.findings-eacl.180, 2023. \url{https://doi.org/10.18653/v1/2023.findings-eacl.180}
\bibitem{ref59} Efficient Transformers: A Survey. doi:10.1145/3530811, 2022. \url{https://doi.org/10.1145/3530811}
\bibitem{ref60} A Survey on Transformer Compression. arXiv:2402.05964, 2024. \url{https://arxiv.org/abs/2402.05964}
\bibitem{ref61} To Transformers and Beyond: Large Language Models for the Genome. arXiv:2311.07621, 2023. \url{https://arxiv.org/abs/2311.07621}
\bibitem{ref62} Attention is All You Need in Speech Separation. arXiv:2010.13154, 2020. \url{https://arxiv.org/abs/2010.13154}
\bibitem{ref63} Overview of the Transformer-based Models for NLP Tasks. doi:10.15439/2020f20, 2020. \url{https://doi.org/10.15439/2020f20}
\bibitem{ref64} On the State of the Art of Evaluation in Neural Language Models. arXiv:1707.05589, 2017. \url{https://arxiv.org/abs/1707.05589}
\bibitem{ref65} N-gram Prediction and Word Difference Representations for Language Modeling. arXiv:2409.03295, 2024. \url{https://arxiv.org/abs/2409.03295}
\bibitem{ref66} Fusing Sentence Embeddings Into LSTM-based Autoregressive Language Models. arXiv:2208.02402, 2022. \url{https://arxiv.org/abs/2208.02402}
\bibitem{ref67} Span Selection Pre-training for Question Answering. doi:10.18653/v1/2020.acl-main.247, 2020. \url{https://doi.org/10.18653/v1/2020.acl-main.247}
\bibitem{ref68} Few-shot Subgoal Planning with Language Models. arXiv:2205.14288, 2022. \url{https://arxiv.org/abs/2205.14288}
\bibitem{ref69} ELECTRA: Pre-training Text Encoders as Discriminators Rather Than Generators. arXiv:2003.10555, 2020. \url{https://arxiv.org/abs/2003.10555}
\bibitem{ref70} Cross-Lingual Supervision improves Large Language Models Pre-training. arXiv:2305.11778, 2023. \url{https://arxiv.org/abs/2305.11778}
\bibitem{ref71} Rho-1: Not All Tokens Are What You Need. arXiv:2404.07965, 2024. \url{https://arxiv.org/abs/2404.07965}
\bibitem{ref72} To Pretrain or Not to Pretrain: Examining the Benefits of Pretrainng on Resource Rich Tasks. doi:10.18653/v1/2020.acl-main.200, 2020. \url{https://doi.org/10.18653/v1/2020.acl-main.200}
\bibitem{ref73} Understanding LLMs: A Comprehensive Overview from Training to Inference. arXiv:2401.02038, 2024. \url{https://arxiv.org/abs/2401.02038}
\bibitem{ref74} Sentence Encoders on STILTs: Supplementary Training on Intermediate Labeled-data Tasks. arXiv:1811.01088, 2018. \url{https://arxiv.org/abs/1811.01088}
\bibitem{ref75} Data Efficient Masked Language Modeling for Vision and Language. arXiv:2109.02040, 2021. \url{https://arxiv.org/abs/2109.02040}
\bibitem{ref76} Hierarchical Multitask Learning Approach for BERT. arXiv:2011.04451, 2020. \url{https://arxiv.org/abs/2011.04451}
\bibitem{ref77} WavLM: Large-Scale Self-Supervised Pre-Training for Full Stack Speech Processing. arXiv:2110.13900, 2021. \url{https://arxiv.org/abs/2110.13900}
\bibitem{ref78} Multimodal Masked Autoencoders Learn Transferable Representations. arXiv:2205.14204, 2022. \url{https://arxiv.org/abs/2205.14204}
\bibitem{ref79} Kimi k1.5: Scaling Reinforcement Learning with LLMs. arXiv:2501.12599, 2025. \url{https://arxiv.org/abs/2501.12599}
\bibitem{ref80} Crafting Efficient Fine-Tuning Strategies for Large Language Models. arXiv:2407.13906, 2024. \url{https://arxiv.org/abs/2407.13906}
\bibitem{ref81} Compacter: Efficient Low-Rank Hypercomplex Adapter Layers. arXiv:2106.04647, 2021. \url{https://arxiv.org/abs/2106.04647}
\bibitem{ref82} Greedy-layer Pruning: Speeding up Transformer Models for Natural Language Processing. arXiv:2105.14839, 2021. \url{https://arxiv.org/abs/2105.14839}
\bibitem{ref83} QEFT: Quantization for Efficient Fine-Tuning of LLMs. arXiv:2410.08661, 2024. \url{https://arxiv.org/abs/2410.08661}
\bibitem{ref84} Self-Distillation Bridges Distribution Gap in Language Model Fine-Tuning. arXiv:2402.13669, 2024. \url{https://arxiv.org/abs/2402.13669}
\bibitem{ref85} MaskPrune: Mask-based LLM Pruning for Layer-wise Uniform Structures. arXiv:2502.14008, 2025. \url{https://arxiv.org/abs/2502.14008}
\bibitem{ref86} ProxSparse: Regularized Learning of Semi-Structured Sparsity Masks for Pretrained LLMs. arXiv:2502.00258, 2025. \url{https://arxiv.org/abs/2502.00258}
\bibitem{ref87} APT: Adaptive Pruning and Tuning Pretrained Language Models for Efficient Training and Inference. arXiv:2401.12200, 2024. \url{https://arxiv.org/abs/2401.12200}
\bibitem{ref88} PruMUX: Augmenting Data Multiplexing with Model Compression. arXiv:2305.14706, 2023. \url{https://arxiv.org/abs/2305.14706}
\bibitem{ref89} Unigram-Normalized Perplexity as a Language Model Performance Measure with Different Vocabulary Sizes. arXiv:2011.13220, 2020. \url{https://arxiv.org/abs/2011.13220}
\bibitem{ref90} GLUE: A Multi-Task Benchmark and Analysis Platform for Natural Language Understanding. doi:10.18653/v1/w18-5446, 2018. \url{https://doi.org/10.18653/v1/w18-5446}
\bibitem{ref91} SuperGLUE: A Stickier Benchmark for General-Purpose Language Understanding Systems. arXiv:1905.00537, 2019. \url{https://arxiv.org/abs/1905.00537}
\bibitem{ref92} CLUE: A Chinese Language Understanding Evaluation Benchmark. arXiv:2004.05986, 2020. \url{https://arxiv.org/abs/2004.05986}
\bibitem{ref93} IndicSUPERB: A Speech Processing Universal Performance Benchmark for Indian languages. arXiv:2208.11761, 2022. \url{https://arxiv.org/abs/2208.11761}
\bibitem{ref94} A global analysis of metrics used for measuring performance in natural language processing. arXiv:2204.11574, 2022. \url{https://arxiv.org/abs/2204.11574}
\bibitem{ref95} Don't Make Your LLM an Evaluation Benchmark Cheater. arXiv:2311.01964, 2023. \url{https://arxiv.org/abs/2311.01964}
\bibitem{ref96} APPLS: Evaluating Evaluation Metrics for Plain Language Summarization. doi:10.18653/v1/2024.emnlp-main.519, 2024. \url{https://doi.org/10.18653/v1/2024.emnlp-main.519}
\bibitem{ref97} BLEU Neighbors: A Reference-less Approach to Automatic Evaluation. doi:10.18653/v1/2020.eval4nlp-1.5, 2020. \url{https://doi.org/10.18653/v1/2020.eval4nlp-1.5}
\bibitem{ref98} Prefix-Tuning: Optimizing Continuous Prompts for Generation. arXiv:2101.00190, 2021. \url{https://arxiv.org/abs/2101.00190}
\bibitem{ref99} Can pruning make Large Language Models more efficient?. arXiv:2310.04573, 2023. \url{https://arxiv.org/abs/2310.04573}
\bibitem{ref100} An exploration of the effect of quantisation on energy consumption and inference time of StarCoder2. arXiv:2411.12758, 2024. \url{https://arxiv.org/abs/2411.12758}
\bibitem{ref101} TED: Accelerate Model Training by Internal Generalization. doi:10.3233/faia240823, 2024. \url{https://doi.org/10.3233/faia240823}
\bibitem{ref102} SPP: Sparsity-Preserved Parameter-Efficient Fine-Tuning for Large Language Models. arXiv:2405.16057, 2024. \url{https://arxiv.org/abs/2405.16057}
\bibitem{ref103} Improving Text Embeddings for Smaller Language Models Using Contrastive Fine-tuning. arXiv:2408.00690, 2024. \url{https://arxiv.org/abs/2408.00690}
\bibitem{ref104} Mini-GPTs: Efficient Large Language Models through Contextual Pruning. arXiv:2312.12682, 2023. \url{https://arxiv.org/abs/2312.12682}
\bibitem{ref105} DistilBERT, a distilled version of BERT: smaller, faster, cheaper and lighter. arXiv:1910.01108, 2019. \url{https://arxiv.org/abs/1910.01108}
\bibitem{ref106} Reinforcement Learning for Aligning Large Language Models Agents with Interactive Environments: Quantifying and Mitigating Prompt Overfitting. arXiv:2410.19920, 2024. \url{https://arxiv.org/abs/2410.19920}
\bibitem{ref107} LLMR: Knowledge Distillation with a Large Language Model-Induced Reward. arXiv:2409.12500, 2024. \url{https://arxiv.org/abs/2409.12500}
\bibitem{ref108} Continuation KD: Improved Knowledge Distillation through the Lens of Continuation Optimization. arXiv:2212.05998, 2022. \url{https://arxiv.org/abs/2212.05998}
\bibitem{ref109} Revisiting Intermediate Layer Distillation for Compressing Language Models: An Overfitting Perspective. arXiv:2302.01530, 2023. \url{https://arxiv.org/abs/2302.01530}
\bibitem{ref110} MixKD: Towards Efficient Distillation of Large-scale Language Models. arXiv:2011.00593, 2020. \url{https://arxiv.org/abs/2011.00593}
\bibitem{ref111} Offline RL for Natural Language Generation with Implicit Language Q Learning. arXiv:2206.11871, 2022. \url{https://arxiv.org/abs/2206.11871}
\bibitem{ref112} Countering Reward Over-optimization in LLM with Demonstration-Guided Reinforcement Learning. arXiv:2404.19409, 2024. \url{https://arxiv.org/abs/2404.19409}
\bibitem{ref113} MEND: Meta dEmonstratioN Distillation for Efficient and Effective In-Context Learning. arXiv:2403.06914, 2024. \url{https://arxiv.org/abs/2403.06914}
\bibitem{ref114} SeCoKD: Aligning Large Language Models for In-Context Learning with Fewer Shots. arXiv:2406.14208, 2024. \url{https://arxiv.org/abs/2406.14208}
\bibitem{ref115} Every Expert Matters: Towards Effective Knowledge Distillation for Mixture-of-Experts Language Models. arXiv:2502.12947, 2025. \url{https://arxiv.org/abs/2502.12947}
\bibitem{ref116} Robust Distillation for Worst-class Performance. arXiv:2206.06479, 2022. \url{https://arxiv.org/abs/2206.06479}
\bibitem{ref117} Cost-effective Distillation of Large Language Models. doi:10.18653/v1/2023.findings-acl.463, 2023. \url{https://doi.org/10.18653/v1/2023.findings-acl.463}
\bibitem{ref118} LoRA-GA: Low-Rank Adaptation with Gradient Approximation. arXiv:2407.05000, 2024. \url{https://arxiv.org/abs/2407.05000}
\bibitem{ref119} LoLDU: Low-Rank Adaptation via Lower-Diag-Upper Decomposition for Parameter-Efficient Fine-Tuning. arXiv:2410.13618, 2024. \url{https://arxiv.org/abs/2410.13618}
\bibitem{ref120} DoRA: Enhancing Parameter-Efficient Fine-Tuning with Dynamic Rank Distribution. arXiv:2405.17357, 2024. \url{https://arxiv.org/abs/2405.17357}
\bibitem{ref121} Sparsity May Be All You Need: Sparse Random Parameter Adaptation. arXiv:2502.15975, 2025. \url{https://arxiv.org/abs/2502.15975}
\bibitem{ref122} SparseAdapter: An Easy Approach for Improving the Parameter-Efficiency of Adapters. arXiv:2210.04284, 2022. \url{https://arxiv.org/abs/2210.04284}
\bibitem{ref123} Parameter-Efficient Fine-Tuning without Introducing New Latency. arXiv:2305.16742, 2023. \url{https://arxiv.org/abs/2305.16742}
\bibitem{ref124} BBTv2: Towards a Gradient-Free Future with Large Language Models. arXiv:2205.11200, 2022. \url{https://arxiv.org/abs/2205.11200}
\bibitem{ref125} LoRA-drop: Efficient LoRA Parameter Pruning based on Output Evaluation. arXiv:2402.07721, 2024. \url{https://arxiv.org/abs/2402.07721}
\bibitem{ref126} PeriodicLoRA: Breaking the Low-Rank Bottleneck in LoRA Optimization. arXiv:2402.16141, 2024. \url{https://arxiv.org/abs/2402.16141}
\bibitem{ref127} CE-LoRA: Computation-Efficient LoRA Fine-Tuning for Language Models. arXiv:2502.01378, 2025. \url{https://arxiv.org/abs/2502.01378}
\bibitem{ref128} SVFT: Parameter-Efficient Fine-Tuning with Singular Vectors. arXiv:2405.19597, 2024. \url{https://arxiv.org/abs/2405.19597}
\bibitem{ref129} AlphaTuning: Quantization-Aware Parameter-Efficient Adaptation of Large-Scale Pre-Trained Language Models. arXiv:2210.03858, 2022. \url{https://arxiv.org/abs/2210.03858}
\bibitem{ref130} X-PEFT: eXtremely Parameter-Efficient Fine-Tuning for Extreme Multi-Profile Scenarios. arXiv:2401.16137, 2024. \url{https://arxiv.org/abs/2401.16137}
\bibitem{ref131} VB-LoRA: Extreme Parameter Efficient Fine-Tuning with Vector Banks. arXiv:2405.15179, 2024. \url{https://arxiv.org/abs/2405.15179}
\bibitem{ref132} Minimal Ranks, Maximum Confidence: Parameter-efficient Uncertainty Quantification for LoRA. arXiv:2502.12122, 2025. \url{https://arxiv.org/abs/2502.12122}
\bibitem{ref133} LoRA Land: 310 Fine-tuned LLMs that Rival GPT-4, A Technical Report. arXiv:2405.00732, 2024. \url{https://arxiv.org/abs/2405.00732}
\bibitem{ref134} PRoLoRA: Partial Rotation Empowers More Parameter-Efficient LoRA. arXiv:2402.16902, 2024. \url{https://arxiv.org/abs/2402.16902}
\bibitem{ref135} GeoLoRA: Geometric integration for parameter efficient fine-tuning. arXiv:2410.18720, 2024. \url{https://arxiv.org/abs/2410.18720}
\bibitem{ref136} ROSA: Random Subspace Adaptation for Efficient Fine-Tuning. arXiv:2407.07802, 2024. \url{https://arxiv.org/abs/2407.07802}
\bibitem{ref137} Self-calibration for Language Model Quantization and Pruning. arXiv:2410.17170, 2024. \url{https://arxiv.org/abs/2410.17170}
\bibitem{ref138} Optimal Brain Compression: A Framework for Accurate Post-Training Quantization and Pruning. arXiv:2208.11580, 2022. \url{https://arxiv.org/abs/2208.11580}
\bibitem{ref139} PD-Quant: Post-Training Quantization Based on Prediction Difference Metric. doi:10.1109/cvpr52729.2023.02340, 2023. \url{https://doi.org/10.1109/cvpr52729.2023.02340}
\bibitem{ref140} Retraining-Based Iterative Weight Quantization for Deep Neural Networks. arXiv:1805.11233, 2018. \url{https://arxiv.org/abs/1805.11233}
\bibitem{ref141} LLM-QAT: Data-Free Quantization Aware Training for Large Language Models. arXiv:2305.17888, 2023. \url{https://arxiv.org/abs/2305.17888}
\bibitem{ref142} QSLR: Post-Training Compression via Quantized Sparse and Low-Rank Factorization. doi:10.1109/access.2025.3615473, 2025. \url{https://doi.org/10.1109/access.2025.3615473}
\bibitem{ref143} Dirichlet Pruning for Neural Network Compression. arXiv:2011.05985, 2020. \url{https://arxiv.org/abs/2011.05985}
\bibitem{ref144} GDP: Stabilized Neural Network Pruning via Gates with Differentiable Polarization. arXiv:2109.02220, 2021. \url{https://arxiv.org/abs/2109.02220}
\bibitem{ref145} Automatic Pruning for Quantized Neural Networks. doi:10.1109/dicta52665.2021.9647074, 2021. \url{https://doi.org/10.1109/dicta52665.2021.9647074}
\bibitem{ref146} DeepTwist: Learning Model Compression via Occasional Weight Distortion. arXiv:1810.12823, 2018. \url{https://arxiv.org/abs/1810.12823}
\bibitem{ref147} Harmonious Coexistence of Structured Weight Pruning and Ternarization for Deep Neural Networks. doi:10.1609/aaai.v34i04.6138, 2020. \url{https://doi.org/10.1609/aaai.v34i04.6138}
\bibitem{ref148} MOHAQ: Multi-Objective Hardware-Aware Quantization of Recurrent Neural Networks. arXiv:2108.01192, 2021. \url{https://arxiv.org/abs/2108.01192}
\bibitem{ref149} VTrans: Accelerating Transformer Compression with Variational Information Bottleneck based Pruning. arXiv:2406.05276, 2024. \url{https://arxiv.org/abs/2406.05276}
\bibitem{ref150} RepQ: Generalizing Quantization-Aware Training for Re-Parametrized Architectures. arXiv:2311.05317, 2023. \url{https://arxiv.org/abs/2311.05317}
\bibitem{ref151} Advances in Pruning and Quantization for Natural Language Processing. doi:10.1109/access.2024.3465631, 2024. \url{https://doi.org/10.1109/access.2024.3465631}
\bibitem{ref152} CRVQ: Channel-relaxed Vector Quantization for Extreme Compression of LLMs. arXiv:2412.09282, 2024. \url{https://arxiv.org/abs/2412.09282}
\bibitem{ref153} Subset-Selection Weight Post-Training Quantization Method for Learned Image Compression Task. doi:10.1109/access.2024.3524553, 2025. \url{https://doi.org/10.1109/access.2024.3524553}
\bibitem{ref154} UltraFeedback Boosting Language Models with High-quality Feedback. arXiv:2310.01377, 2023. \url{https://arxiv.org/abs/2310.01377}
\bibitem{ref155} The Atari Grand Challenge Dataset. arXiv:1705.10998, 2017. \url{https://arxiv.org/abs/1705.10998}
\bibitem{ref156} On the Challenges and Practices of Reinforcement Learning from Real Human Feedback. doi:10.1007/978-3-031-74627-7\_21, 2024. \url{https://doi.org/10.1007/978-3-031-74627-7_21}
\bibitem{ref157} Iterative Data Smoothing: Mitigating Reward Overfitting and Overoptimization in RLHF. arXiv:2401.16335, 2024. \url{https://arxiv.org/abs/2401.16335}
\bibitem{ref158} RLAIF Scaling Reinforcement Learning from Human Feedback with AI Feedback. arXiv:2309.00267, 2023. \url{https://arxiv.org/abs/2309.00267}
\bibitem{ref159} MaxMin-RLHF Towards Equitable Alignment of Large Language Models with Diverse Human Preferences. arXiv:2402.08925, 2024. \url{https://arxiv.org/abs/2402.08925}
\bibitem{ref160} Reinforcement Learning in Healthcare: A Survey. arXiv:1908.08796, 2019. \url{https://arxiv.org/abs/1908.08796}
\bibitem{ref161} System Design Perspective for Human-Level Agents Using Deep Reinforcement Learning: A Survey. doi:10.1109/access.2017.2777827, 2017. \url{https://doi.org/10.1109/access.2017.2777827}
\bibitem{ref162} Reinforcement Learning Algorithms and Applications in Healthcare and Robotics: A Comprehensive and Systematic Review. doi:10.3390/s24082461, 2024. \url{https://doi.org/10.3390/s24082461}
\bibitem{ref163} Uni-RLHF: Universal Platform and Benchmark Suite for Reinforcement Learning with Diverse Human Feedback. arXiv:2402.02423, 2024. \url{https://arxiv.org/abs/2402.02423}
\bibitem{ref164} Reinforcement Learning Enhanced LLMs: A Survey. arXiv:2412.10400, 2024. \url{https://arxiv.org/abs/2412.10400}
\bibitem{ref165} Retroformer: Retrospective Large Language Agents with Policy Gradient Optimization. arXiv:2308.02151, 2023. \url{https://arxiv.org/abs/2308.02151}
\bibitem{ref166} On the Theory of Policy Gradient Methods: Optimality, Approximation, and Distribution Shift. arXiv:1908.00261, 2019. \url{https://arxiv.org/abs/1908.00261}
\bibitem{ref167} Policy Learning with a Natural Language Action Space: A Causal Approach. arXiv:2502.17538, 2025. \url{https://arxiv.org/abs/2502.17538}
\bibitem{ref168} How Can LLM Guide RL? A Value-Based Approach. arXiv:2402.16181, 2024. \url{https://arxiv.org/abs/2402.16181}
\bibitem{ref169} Implementation Matters in Deep Policy Gradients: A Case Study on PPO and TRPO. arXiv:2005.12729, 2020. \url{https://arxiv.org/abs/2005.12729}
\bibitem{ref170} ArCHer: Training Language Model Agents via Hierarchical Multi-Turn RL. arXiv:2402.19446, 2024. \url{https://arxiv.org/abs/2402.19446}
\bibitem{ref171} FLAG-Trader: Fusion LLM-Agent with Gradient-based Reinforcement Learning for Financial Trading. arXiv:2502.11433, 2025. \url{https://arxiv.org/abs/2502.11433}
\bibitem{ref172} Semi-On-Policy Training for Sample Efficient Multi-Agent Policy Gradients. arXiv:2104.13446, 2021. \url{https://arxiv.org/abs/2104.13446}
\bibitem{ref173} Optima: Optimizing Effectiveness and Efficiency for LLM-Based Multi-Agent System. arXiv:2410.08115, 2024. \url{https://arxiv.org/abs/2410.08115}
\bibitem{ref174} COMMA: A Communicative Multimodal Multi-Agent Benchmark. arXiv:2410.07553, 2024. \url{https://arxiv.org/abs/2410.07553}
\bibitem{ref175} Cooperative Pathfinding based on Multi-agent RRT* Fixed Node. arXiv:1911.07840, 2019. \url{https://arxiv.org/abs/1911.07840}
\bibitem{ref176} AgentVerse: Facilitating Multi-Agent Collaboration and Exploring Emergent Behaviors in Agents. arXiv:2308.10848, 2023. \url{https://arxiv.org/abs/2308.10848}
\bibitem{ref177} Resource allocation in dynamic multiagent systems. arXiv:2102.08317, 2021. \url{https://arxiv.org/abs/2102.08317}
\bibitem{ref178} Melting Pot 2.0. arXiv:2211.13746, 2022. \url{https://arxiv.org/abs/2211.13746}
\bibitem{ref179} Mixed-Initiative Human-Robot Teaming under Suboptimality with Online Bayesian Adaptation. arXiv:2403.16178, 2024. \url{https://arxiv.org/abs/2403.16178}
\bibitem{ref180} Optimization under Attack: Resilience, Vulnerability, and the Path to Collapse. arXiv:2502.05954, 2025. \url{https://arxiv.org/abs/2502.05954}
\bibitem{ref181} MACRPO: Multi-agent cooperative recurrent policy optimization. doi:10.3389/frobt.2024.1394209, 2024. \url{https://doi.org/10.3389/frobt.2024.1394209}
\bibitem{ref182} Interaction Pattern Disentangling for Multi-Agent Reinforcement Learning. doi:10.1109/tpami.2024.3399936, 2024. \url{https://doi.org/10.1109/tpami.2024.3399936}
\bibitem{ref183} Scaling Submodular Optimization Approaches for Control Applications in Networked Systems. arXiv:1810.02837, 2018. \url{https://arxiv.org/abs/1810.02837}
\bibitem{ref184} Towards a Standardised Performance Evaluation Protocol for Cooperative MARL. arXiv:2209.10485, 2022. \url{https://arxiv.org/abs/2209.10485}
\bibitem{ref185} An Ideal Team Is More than a Team of Ideal Agents. doi:10.3233/faia200074, 2020. \url{https://doi.org/10.3233/faia200074}
\bibitem{ref186} A Tutorial on the Structure of Distributed Optimization Algorithms. doi:10.1109/cdc49753.2023.10383319, 2023. \url{https://doi.org/10.1109/cdc49753.2023.10383319}
\bibitem{ref187} Shades of BLEU, Flavours of Success: The Case of MultiWOZ. doi:10.18653/v1/2021.gem-1.4, 2021. \url{https://doi.org/10.18653/v1/2021.gem-1.4}
\bibitem{ref188} AgentQuest: A Modular Benchmark Framework to Measure Progress and Improve LLM Agents. arXiv:2404.06411, 2024. \url{https://arxiv.org/abs/2404.06411}
\bibitem{ref189} MetricsVis: A Visual Analytics System for Evaluating Employee Performance in Public Safety Agencies. doi:10.1109/tvcg.2019.2934603, 2019. \url{https://doi.org/10.1109/tvcg.2019.2934603}
\bibitem{ref190} Efficiency Metrics for Data-Driven Models: A Text Summarization Case Study. doi:10.18653/v1/w19-8630, 2019. \url{https://doi.org/10.18653/v1/w19-8630}
\bibitem{ref191} Constructing and meta-evaluating state-aware evaluation metrics for interactive search systems. doi:10.1007/s10791-023-09426-1, 2023. \url{https://doi.org/10.1007/s10791-023-09426-1}
\bibitem{ref192} Learning from Task Descriptions. doi:10.18653/v1/2020.emnlp-main.105, 2020. \url{https://doi.org/10.18653/v1/2020.emnlp-main.105}
\bibitem{ref193} BEAMetrics: A Benchmark for Language Generation Evaluation Evaluation. arXiv:2110.09147, 2021. \url{https://arxiv.org/abs/2110.09147}
\bibitem{ref194} FineD-Eval: Fine-grained Automatic Dialogue-Level Evaluation. arXiv:2210.13832, 2022. \url{https://arxiv.org/abs/2210.13832}
\bibitem{ref195} CTRLEval: An Unsupervised Reference-Free Metric for Evaluating Controlled Text Generation. arXiv:2204.00862, 2022. \url{https://arxiv.org/abs/2204.00862}
\bibitem{ref196} VLUE: A Multi-Task Benchmark for Evaluating Vision-Language Models. arXiv:2205.15237, 2022. \url{https://arxiv.org/abs/2205.15237}
\bibitem{ref197} SOTVerse: A User-Defined Task Space of Single Object Tracking. doi:10.1007/s11263-023-01908-5, 2023. \url{https://doi.org/10.1007/s11263-023-01908-5}
\bibitem{ref198} BALI—A Benchmark for Accelerated Language Model Inference. doi:10.1109/access.2025.3576898, 2025. \url{https://doi.org/10.1109/access.2025.3576898}
\bibitem{ref199} P-MMEval: A Parallel Multilingual Multitask Benchmark for Consistent Evaluation of LLMs. arXiv:2411.09116, 2024. \url{https://arxiv.org/abs/2411.09116}
\bibitem{ref200} M5 - A Diverse Benchmark to Assess the Performance of Large Multimodal Models Across Multilingual and Multicultural Vision-Language Tasks. arXiv:2407.03791, 2024. \url{https://arxiv.org/abs/2407.03791}
\bibitem{ref201} LiveCodeBench Holistic and Contamination Free Evaluation of Large Language Models for Code. arXiv:2403.07974, 2024. \url{https://arxiv.org/abs/2403.07974}
\bibitem{ref202} MMBench: Is Your Multi-modal Model an All-around Player?. arXiv:2307.06281, 2023. \url{https://arxiv.org/abs/2307.06281}
\bibitem{ref203} Evaluating Human-AI Collaboration: A Review and Methodological Framework. arXiv:2407.19098, 2024. \url{https://arxiv.org/abs/2407.19098}
\bibitem{ref204} The Harmony Index: a Utilitarian Metric for Measuring Effectiveness in Mixed-Skill Teams. arXiv:2203.12222, 2022. \url{https://arxiv.org/abs/2203.12222}
\bibitem{ref205} User experience testing methods: Conclusions from the literature. doi:10.55214/25768484.v8i5.1843, 2024. \url{https://doi.org/10.55214/25768484.v8i5.1843}
\bibitem{ref206} User-Centric Evaluation of Recommender Systems in Social Learning Platforms: Accuracy is Just the Tip of the Iceberg. doi:10.1109/tlt.2017.2732349, 2017. \url{https://doi.org/10.1109/tlt.2017.2732349}
\bibitem{ref207} Explainable Artificial Intelligence: Evaluating the Objective and Subjective Impacts of xAI on Human-Agent Interaction. doi:10.1080/10447318.2022.2101698, 2022. \url{https://doi.org/10.1080/10447318.2022.2101698}
\bibitem{ref208} Investigations of Performance and Bias in Human-AI Teamwork in Hiring. arXiv:2202.11812, 2022. \url{https://arxiv.org/abs/2202.11812}
\bibitem{ref209} Technology-Enhanced Learning: An Optimal CPS Learning Application. doi:10.3390/su11164415, 2019. \url{https://doi.org/10.3390/su11164415}
\bibitem{ref210} Human-Centric Foundation Models: Perception, Generation and Agentic Modeling. arXiv:2502.08556, 2025. \url{https://arxiv.org/abs/2502.08556}
\bibitem{ref211} Effective Task Allocation in Ad Hoc Human-Agent Teams. doi:10.3233/faia220197, 2022. \url{https://doi.org/10.3233/faia220197}
\bibitem{ref212} Two Many Cooks: Understanding Dynamic Human-Agent Team Communication and Perception Using Overcooked 2. arXiv:2110.03071, 2021. \url{https://arxiv.org/abs/2110.03071}
\bibitem{ref213} Human-Centric Research for NLP: Towards a Definition and Guiding Questions. arXiv:2207.04447, 2022. \url{https://arxiv.org/abs/2207.04447}
\bibitem{ref214} Interactive Speculative Planning: Enhance Agent Efficiency through Co-design of System and User Interface. arXiv:2410.00079, 2024. \url{https://arxiv.org/abs/2410.00079}
\bibitem{ref215} Evaluating Mixed and Augmented Reality: A Systematic Literature Review (2009-2019). doi:10.1109/ismar50242.2020.00069, 2020. \url{https://doi.org/10.1109/ismar50242.2020.00069}
\bibitem{ref216} Persona preparedness: a survey instrument for measuring the organizational readiness for deploying personas. doi:10.1007/s10799-022-00373-9, 2022. \url{https://doi.org/10.1007/s10799-022-00373-9}
\bibitem{ref217} CREW: Facilitating Human-AI Teaming Research. arXiv:2408.00170, 2024. \url{https://arxiv.org/abs/2408.00170}
\bibitem{ref218} Towards Objective Evaluation of Socially-Situated Conversational Robots: Assessing Human-Likeness through Multimodal User Behaviors. doi:10.1145/3610661.3617151, 2023. \url{https://doi.org/10.1145/3610661.3617151}
\bibitem{ref219} A Comprehensive Analysis of a Social Intelligence Dataset and Response Tendencies Between Large Language Models (LLMs) and Humans. doi:10.3390/s25020477, 2025. \url{https://doi.org/10.3390/s25020477}
\bibitem{ref220} Is my Meeting Summary Good? Estimating Quality with a Multi-LLM Evaluator. arXiv:2411.18444, 2024. \url{https://arxiv.org/abs/2411.18444}
\bibitem{ref221} The role of socio-emotional attributes in enhancing human-AI collaboration. doi:10.3389/fpsyg.2024.1369957, 2024. \url{https://doi.org/10.3389/fpsyg.2024.1369957}
\bibitem{ref222} Optimization Techniques for Sentiment Analysis Based on LLM (GPT-3). doi:10.54254/2755-2721/67/2024ma0060, 2024. \url{https://doi.org/10.54254/2755-2721/67/2024ma0060}
\bibitem{ref223} Large Scale Fine-Tuned Transformers Models Application for Business Names Generation. doi:10.31577/cai\_2023\_3\_525, 2023. \url{https://doi.org/10.31577/cai_2023_3_525}
\bibitem{ref224} Investigating translation for Indic languages with BLOOMZ-3b through prompting and LoRA fine-tuning. doi:10.1038/s41598-024-74617-9, 2024. \url{https://doi.org/10.1038/s41598-024-74617-9}
\bibitem{ref225} Sentiment Analysis in the Era of Large Language Models: A Reality Check. arXiv:2305.15005, 2023. \url{https://arxiv.org/abs/2305.15005}
\bibitem{ref226} The Model Arena for Cross-lingual Sentiment Analysis: A Comparative Study in the Era of Large Language Models. arXiv:2406.19358, 2024. \url{https://arxiv.org/abs/2406.19358}
\bibitem{ref227} Overcoming Language Disparity in Online Content Classification with Multimodal Learning. arXiv:2205.09744, 2022. \url{https://arxiv.org/abs/2205.09744}
\bibitem{ref228} The Ultimate Guide to Fine-Tuning LLMs from Basics to Breakthroughs: An Exhaustive Review of Technologies, Research, Best Practices, Applied Research Challenges and Opportunities. arXiv:2408.13296, 2024. \url{https://arxiv.org/abs/2408.13296}
\bibitem{ref229} Transformers in Vision: A Survey. doi:10.1145/3505244, 2022. \url{https://doi.org/10.1145/3505244}
\bibitem{ref230} A Review of Transformer-Based Models for Computer Vision Tasks: Capturing Global Context and Spatial Relationships. arXiv:2408.15178, 2024. \url{https://arxiv.org/abs/2408.15178}
\bibitem{ref231} Toward Transformer-Based Object Detection. arXiv:2012.09958, 2020. \url{https://arxiv.org/abs/2012.09958}
\bibitem{ref232} YOLO-RD: Introducing Relevant and Compact Explicit Knowledge to YOLO by Retriever-Dictionary. arXiv:2410.15346, 2024. \url{https://arxiv.org/abs/2410.15346}
\bibitem{ref233} MuLan Multimodal-LLM Agent for Progressive Multi-Object Diffusion. arXiv:2402.12741, 2024. \url{https://arxiv.org/abs/2402.12741}
\bibitem{ref234} Uni-Perceiver v2: A Generalist Model for Large-Scale Vision and Vision-Language Tasks. doi:10.1109/cvpr52729.2023.00264, 2023. \url{https://doi.org/10.1109/cvpr52729.2023.00264}
\bibitem{ref235} mPLUG: Effective and Efficient Vision-Language Learning by Cross-modal Skip-connections. arXiv:2205.12005, 2022. \url{https://arxiv.org/abs/2205.12005}
\bibitem{ref236} Grounding Language Models in Autonomous Loco-manipulation Tasks. arXiv:2409.01326, 2024. \url{https://arxiv.org/abs/2409.01326}
\bibitem{ref237} LLM-Based Human-Robot Collaboration Framework for Manipulation Tasks. arXiv:2308.14972, 2023. \url{https://arxiv.org/abs/2308.14972}
\bibitem{ref238} Leveraging large language models for autonomous robotic mapping and navigation. doi:10.1177/17298806251325965, 2025. \url{https://doi.org/10.1177/17298806251325965}
\bibitem{ref239} Speech-Guided Sequential Planning for Autonomous Navigation using Large Language Model Meta AI 3 (Llama3). arXiv:2407.09890, 2024. \url{https://arxiv.org/abs/2407.09890}
\bibitem{ref240} The Conversation is the Command: Interacting with Real-World Autonomous Robots Through Natural Language. doi:10.1145/3610978.3640723, 2024. \url{https://doi.org/10.1145/3610978.3640723}
\bibitem{ref241} Multimodal Human-Autonomous Agents Interaction Using Pre-Trained Language and Visual Foundation Models. arXiv:2403.12273, 2024. \url{https://arxiv.org/abs/2403.12273}
\bibitem{ref242} RoboCoder: Robotic Learning from Basic Skills to General Tasks with Large Language Models. arXiv:2406.03757, 2024. \url{https://arxiv.org/abs/2406.03757}
\bibitem{ref243} REVERIE: Remote Embodied Visual Referring Expression in Real Indoor Environments. doi:10.1109/cvpr42600.2020.01000, 2020. \url{https://doi.org/10.1109/cvpr42600.2020.01000}
\bibitem{ref244} CAMON: Cooperative Agents for Multi-Object Navigation with LLM-based Conversations. arXiv:2407.00632, 2024. \url{https://arxiv.org/abs/2407.00632}
\bibitem{ref245} How can I help you? An Intelligent Virtual Assistant for Industrial Robots. doi:10.1145/3434074.3447163, 2021. \url{https://doi.org/10.1145/3434074.3447163}
\bibitem{ref246} Stars, Stripes, and Silicon: Unravelling the ChatGPT's All-American, Monochrome, Cis-centric Bias. arXiv:2410.13868, 2024. \url{https://arxiv.org/abs/2410.13868}
\bibitem{ref247} Double Jeopardy and Climate Impact in the Use of Large Language Models: Socio-economic Disparities and Reduced Utility for Non-English Speakers. arXiv:2410.10665, 2024. \url{https://arxiv.org/abs/2410.10665}
\bibitem{ref248} BiasDPO: Mitigating Bias in Language Models through Direct Preference Optimization. arXiv:2407.13928, 2024. \url{https://arxiv.org/abs/2407.13928}
\bibitem{ref249} Towards Resource Efficient and Interpretable Bias Mitigation in Large Language Models. arXiv:2412.01711, 2024. \url{https://arxiv.org/abs/2412.01711}
\bibitem{ref250} Data Augmentation for Neural NLP. arXiv:2302.11412, 2023. \url{https://arxiv.org/abs/2302.11412}
\bibitem{ref251} Increasing Diversity While Maintaining Accuracy: Text Data Generation with Large Language Models and Human Interventions. doi:10.18653/v1/2023.acl-long.34, 2023. \url{https://doi.org/10.18653/v1/2023.acl-long.34}
\bibitem{ref252} Cultural Bias in Large Language Models: A Comprehensive Analysis and Mitigation Strategies. doi:10.1515/jtc-2023-0019, 2023. \url{https://doi.org/10.1515/jtc-2023-0019}
\bibitem{ref253} Distributed Inference Performance Optimization for LLMs on CPUs. arXiv:2407.00029, 2024. \url{https://arxiv.org/abs/2407.00029}
\bibitem{ref254} Towards Fast Multilingual LLM Inference: Speculative Decoding and Specialized Drafters. arXiv:2406.16758, 2024. \url{https://arxiv.org/abs/2406.16758}
\bibitem{ref255} SplitLLM: Collaborative Inference of LLMs for Model Placement and Throughput Optimization. arXiv:2410.10759, 2024. \url{https://arxiv.org/abs/2410.10759}
\bibitem{ref256} Scaling Laws for Upcycling Mixture-of-Experts Language Models. arXiv:2502.03009, 2025. \url{https://arxiv.org/abs/2502.03009}
\bibitem{ref257} Scaling Retrieval-Based Language Models with a Trillion-Token Datastore. arXiv:2407.12854, 2024. \url{https://arxiv.org/abs/2407.12854}
\bibitem{ref258} Efficient and scalable reinforcement learning for large-scale network control. doi:10.1038/s42256-024-00879-7, 2024. \url{https://doi.org/10.1038/s42256-024-00879-7}
\bibitem{ref259} Towards Optimal Caching and Model Selection for Large Model Inference. doi:10.52202/075280-2577, 2023. \url{https://doi.org/10.52202/075280-2577}
\bibitem{ref260} GLaM: Efficient Scaling of Language Models with Mixture-of-Experts. arXiv:2112.06905, 2021. \url{https://arxiv.org/abs/2112.06905}
\bibitem{ref261} SMoA: Improving Multi-agent Large Language Models with Sparse Mixture-of-Agents. arXiv:2411.03284, 2024. \url{https://arxiv.org/abs/2411.03284}
\bibitem{ref262} Understanding the Weakness of Large Language Model Agents within a Complex Android Environment. doi:10.1145/3637528.3671650, 2024. \url{https://doi.org/10.1145/3637528.3671650}
\bibitem{ref263} Breaking Language Barriers: Cross-Lingual Continual Pre-Training at Scale. arXiv:2407.02118, 2024. \url{https://arxiv.org/abs/2407.02118}
\bibitem{ref264} InfLLM Unveiling the Intrinsic Capacity of LLMs for Understanding Extremely Long Sequences with Training-Free Memory. arXiv:2402.04617, 2024. \url{https://arxiv.org/abs/2402.04617}
\bibitem{ref265} Explainability for Large Language Models: A Survey. doi:10.1145/3639372, 2024. \url{https://doi.org/10.1145/3639372}
\bibitem{ref266} Mechanistic interpretability of large language models with applications to the financial services industry. doi:10.1145/3677052.3698612, 2024. \url{https://doi.org/10.1145/3677052.3698612}
\bibitem{ref267} Interpreting token compositionality in LLMs: A robustness analysis. arXiv:2410.12924, 2024. \url{https://arxiv.org/abs/2410.12924}
\bibitem{ref268} Understanding the Effect of Algorithm Transparency of Model Explanations in Text-to-SQL Semantic Parsing. arXiv:2410.16283, 2024. \url{https://arxiv.org/abs/2410.16283}
\bibitem{ref269} Explainable and Interpretable Multimodal Large Language Models: A Comprehensive Survey. arXiv:2412.02104, 2024. \url{https://arxiv.org/abs/2412.02104}
\bibitem{ref270} TokenSHAP: Interpreting Large Language Models with Monte Carlo Shapley Value Estimation. arXiv:2407.10114, 2024. \url{https://arxiv.org/abs/2407.10114}
\bibitem{ref271} Contrastive Explanations for Reinforcement Learning in terms of Expected Consequences. arXiv:1807.08706, 2018. \url{https://arxiv.org/abs/1807.08706}
\bibitem{ref272} From large language models to small logic programs: building global explanations from disagreeing local post-hoc explainers. doi:10.1007/s10458-024-09663-8, 2024. \url{https://doi.org/10.1007/s10458-024-09663-8}
\bibitem{ref273} Bridging AI and Human Understanding: Interpretable Deep Learning in Practice. doi:10.52783/jier.v4i3.2200, 2024. \url{https://doi.org/10.52783/jier.v4i3.2200}
\bibitem{ref274} VITA-1.5: Towards GPT-4o Level Real-Time Vision and Speech Interaction. arXiv:2501.01957, 2025. \url{https://arxiv.org/abs/2501.01957}
\bibitem{ref275} 4M: Massively Multimodal Masked Modeling. arXiv:2312.06647, 2023. \url{https://arxiv.org/abs/2312.06647}
\bibitem{ref276} MM1.5: Methods, Analysis \& Insights from Multimodal LLM Fine-tuning. arXiv:2409.20566, 2024. \url{https://arxiv.org/abs/2409.20566}
\bibitem{ref277} Valley: Video Assistant with Large Language Model Enhanced Ability. doi:10.1145/3796716, 2023. \url{https://doi.org/10.1145/3796716}
\bibitem{ref278} Discrete Multimodal Transformers with a Pretrained Large Language Model for Mixed-Supervision Speech Processing. arXiv:2406.06582, 2024. \url{https://arxiv.org/abs/2406.06582}
\bibitem{ref279} Large-scale Multi-modal Pre-trained Models: A Comprehensive Survey. doi:10.1007/s11633-022-1410-8, 2023. \url{https://doi.org/10.1007/s11633-022-1410-8}
\bibitem{ref280} Mixture-of-Transformers: A Sparse and Scalable Architecture for Multi-Modal Foundation Models. arXiv:2411.04996, 2024. \url{https://arxiv.org/abs/2411.04996}
\bibitem{ref281} FLAVA: A Foundational Language And Vision Alignment Model. doi:10.1109/cvpr52688.2022.01519, 2022. \url{https://doi.org/10.1109/cvpr52688.2022.01519}
\bibitem{ref282} S3: A Simple Strong Sample-effective Multimodal Dialog System. arXiv:2406.18305, 2024. \url{https://arxiv.org/abs/2406.18305}
\bibitem{ref283} Stress Detection using Multimodal Representation Learning, Fusion Techniques, and Applications. doi:10.52783/jisem.v10i16s.2593, 2025. \url{https://doi.org/10.52783/jisem.v10i16s.2593}
\bibitem{ref284} MammothModa: Multi-Modal Large Language Model. arXiv:2406.18193, 2024. \url{https://arxiv.org/abs/2406.18193}
\bibitem{ref285} Uni-MoE: Scaling Unified Multimodal LLMs With Mixture of Experts. doi:10.1109/tpami.2025.3532688, 2025. \url{https://doi.org/10.1109/tpami.2025.3532688}
\bibitem{ref286} Mr. Right: Multimodal Retrieval on Representation of ImaGe witH Text. arXiv:2209.13764, 2022. \url{https://arxiv.org/abs/2209.13764}
\bibitem{ref287} What Large Language Models Bring to Text-rich VQA?. arXiv:2311.07306, 2023. \url{https://arxiv.org/abs/2311.07306}
\bibitem{ref288} Few-Shot VQA with Frozen LLMs: A Tale of Two Approaches. arXiv:2403.11317, 2024. \url{https://arxiv.org/abs/2403.11317}
\bibitem{ref289} FineCaption: Compositional Image Captioning Focusing on Wherever You Want at Any Granularity. arXiv:2411.15411, 2024. \url{https://arxiv.org/abs/2411.15411}
\bibitem{ref290} Contextual Object Detection with Multimodal Large Language Models. doi:10.1007/s11263-024-02214-4, 2024. \url{https://doi.org/10.1007/s11263-024-02214-4}
\bibitem{ref291} Sheffield MultiMT: Using Object Posterior Predictions for Multimodal Machine Translation. doi:10.18653/v1/w17-4752, 2017. \url{https://doi.org/10.18653/v1/w17-4752}
\bibitem{ref292} VAST A Vision-Audio-Subtitle-Text Omni-Modality Foundation Model and Dataset. arXiv:2305.18500, 2023. \url{https://arxiv.org/abs/2305.18500}
\bibitem{ref293} 12-in-1: Multi-Task Vision and Language Representation Learning. doi:10.1109/cvpr42600.2020.01045, 2020. \url{https://doi.org/10.1109/cvpr42600.2020.01045}
\bibitem{ref294} ICU: Conquering Language Barriers in Vision-and-Language Modeling by Dividing the Tasks into Image Captioning and Language Understanding. arXiv:2310.12531, 2023. \url{https://arxiv.org/abs/2310.12531}
\bibitem{ref295} PaLI-X: On Scaling up a Multilingual Vision and Language Model. arXiv:2305.18565, 2023. \url{https://arxiv.org/abs/2305.18565}
\bibitem{ref296} A survey on advancements in image-text multimodal models: From general techniques to biomedical implementations. doi:10.1016/j.compbiomed.2024.108709, 2024. \url{https://doi.org/10.1016/j.compbiomed.2024.108709}
\bibitem{ref297} A survey on multimodal large language models. doi:10.1093/nsr/nwae403, 2023. \url{https://doi.org/10.1093/nsr/nwae403}
\bibitem{ref298} Multimodal Machine Learning: A Survey and Taxonomy. arXiv:1705.09406, 2017. \url{https://arxiv.org/abs/1705.09406}
\bibitem{ref299} The Troubling Emergence of Hallucination in Large Language Models - An Extensive Definition, Quantification, and Prescriptive Remediations. arXiv:2310.04988, 2023. \url{https://arxiv.org/abs/2310.04988}
\bibitem{ref300} MIND: Modality-Informed Knowledge Distillation Framework for Multimodal Clinical Prediction Tasks. arXiv:2502.01158, 2025. \url{https://arxiv.org/abs/2502.01158}
\bibitem{ref301} A Review on Methods and Applications in Multimodal Deep Learning. doi:10.1145/3545572, 2022. \url{https://doi.org/10.1145/3545572}
\bibitem{ref302} Detecting and Evaluating Medical Hallucinations in Large Vision Language Models. arXiv:2406.10185, 2024. \url{https://arxiv.org/abs/2406.10185}
\bibitem{ref303} VALOR-EVAL: Holistic Coverage and Faithfulness Evaluation of Large Vision-Language Models. arXiv:2404.13874, 2024. \url{https://arxiv.org/abs/2404.13874}
\bibitem{ref304} From Efficient Multimodal Models to World Models: A Survey. arXiv:2407.00118, 2024. \url{https://arxiv.org/abs/2407.00118}
\bibitem{ref305} NoiseBoost: Alleviating Hallucination with Noise Perturbation for Multimodal Large Language Models. arXiv:2405.20081, 2024. \url{https://arxiv.org/abs/2405.20081}
\bibitem{ref306} CentralNet: a Multilayer Approach for Multimodal Fusion. arXiv:1808.07275, 2018. \url{https://arxiv.org/abs/1808.07275}
\bibitem{ref307} A Survey on Hallucination in Large Language Models: Principles, Taxonomy, Challenges, and Open Questions. doi:10.1145/3703155, 2024. \url{https://doi.org/10.1145/3703155}
\bibitem{ref308} Bias and Fairness in Large Language Models: A Survey. doi:10.1162/coli\_a\_00524, 2024. \url{https://doi.org/10.1162/coli_a_00524}
\bibitem{ref309} Towards Understanding and Mitigating Social Biases in Language Models. arXiv:2106.13219, 2021. \url{https://arxiv.org/abs/2106.13219}
\bibitem{ref310} Locating and Mitigating Gender Bias in Large Language Models. arXiv:2403.14409, 2024. \url{https://arxiv.org/abs/2403.14409}
\bibitem{ref311} Fairness-Aware Structured Pruning in Transformers. arXiv:2312.15398, 2023. \url{https://arxiv.org/abs/2312.15398}
\bibitem{ref312} Unveiling and Mitigating Bias in Mental Health Analysis with Large Language Models. arXiv:2406.12033, 2024. \url{https://arxiv.org/abs/2406.12033}
\bibitem{ref313} Rejected Dialects: Biases Against African American Language in Reward Models. arXiv:2502.12858, 2025. \url{https://arxiv.org/abs/2502.12858}
\bibitem{ref314} Mitigating Language-Dependent Ethnic Bias in BERT. doi:10.18653/v1/2021.emnlp-main.42, 2021. \url{https://doi.org/10.18653/v1/2021.emnlp-main.42}
\bibitem{ref315} Toward Fairness in Text Generation via Mutual Information Minimization based on Importance Sampling. arXiv:2302.13136, 2023. \url{https://arxiv.org/abs/2302.13136}
\bibitem{ref316} FairDistillation: Mitigating Stereotyping in Language Models. arXiv:2207.04546, 2022. \url{https://arxiv.org/abs/2207.04546}
\bibitem{ref317} Indian-BhED: A Dataset for Measuring India-Centric Biases in Large Language Models. doi:10.1145/3677525.3678666, 2024. \url{https://doi.org/10.1145/3677525.3678666}
\bibitem{ref318} Unmasking Implicit Bias: Evaluating Persona-Prompted LLM Responses in Power-Disparate Social Scenarios. doi:10.18653/v1/2025.naacl-long.50, 2025. \url{https://doi.org/10.18653/v1/2025.naacl-long.50}
\bibitem{ref319} Social Bias Probing: Fairness Benchmarking for Language Models. doi:10.18653/v1/2024.emnlp-main.812, 2024. \url{https://doi.org/10.18653/v1/2024.emnlp-main.812}
\bibitem{ref320} The Impossibility of Fair LLMs. arXiv:2406.03198, 2024. \url{https://arxiv.org/abs/2406.03198}
\bibitem{ref321} Diversity and language technology: how language modeling bias causes epistemic injustice. doi:10.1007/s10676-023-09742-6, 2024. \url{https://doi.org/10.1007/s10676-023-09742-6}
\bibitem{ref322} FrenchToxicityPrompts: a Large Benchmark for Evaluating and Mitigating Toxicity in French Texts. arXiv:2406.17566, 2024. \url{https://arxiv.org/abs/2406.17566}
\bibitem{ref323} Safe to Serve: Aligning Instruction-Tuned Models for Safety and Helpfulness. arXiv:2412.00074, 2024. \url{https://arxiv.org/abs/2412.00074}
\bibitem{ref324} Can LLMs Rank the Harmfulness of Smaller LLMs? We are Not There Yet. arXiv:2502.05291, 2025. \url{https://arxiv.org/abs/2502.05291}
\bibitem{ref325} Self-Detoxifying Language Models via Toxification Reversal. arXiv:2310.09573, 2023. \url{https://arxiv.org/abs/2310.09573}
\bibitem{ref326} SafetyPrompts a Systematic Review of Open Datasets for Evaluating and Improving Large Language Model Safety. arXiv:2404.05399, 2024. \url{https://arxiv.org/abs/2404.05399}
\bibitem{ref327} ALERT: A Comprehensive Benchmark for Assessing Large Language Models' Safety through Red Teaming. arXiv:2404.08676, 2024. \url{https://arxiv.org/abs/2404.08676}
\bibitem{ref328} SimpleSafetyTests: a Test Suite for Identifying Critical Safety Risks in Large Language Models. arXiv:2311.08370, 2023. \url{https://arxiv.org/abs/2311.08370}
\bibitem{ref329} Do-Not-Answer: A Dataset for Evaluating Safeguards in LLMs. arXiv:2308.13387, 2023. \url{https://arxiv.org/abs/2308.13387}
\bibitem{ref330} Fine-Grained Detoxification via Instance-Level Prefixes for Large Language Models. arXiv:2402.15202, 2024. \url{https://arxiv.org/abs/2402.15202}
\bibitem{ref331} PARDEN, Can You Repeat That? Defending against Jailbreaks via Repetition. arXiv:2405.07932, 2024. \url{https://arxiv.org/abs/2405.07932}
\bibitem{ref332} Risk-Averse Fine-tuning of Large Language Models. arXiv:2501.06911, 2025. \url{https://arxiv.org/abs/2501.06911}
\bibitem{ref333} ToViLaG: Your Visual-Language Generative Model is Also An Evildoer. doi:10.18653/v1/2023.emnlp-main.213, 2023. \url{https://doi.org/10.18653/v1/2023.emnlp-main.213}
\bibitem{ref334} Contrastive Perplexity for Controlled Generation: An Application in Detoxifying Large Language Models. arXiv:2401.08491, 2024. \url{https://arxiv.org/abs/2401.08491}
\bibitem{ref335} LLMGuard: Guarding against Unsafe LLM Behavior. doi:10.1609/aaai.v38i21.30566, 2024. \url{https://doi.org/10.1609/aaai.v38i21.30566}
\bibitem{ref336} Language Generation Models Can Cause Harm: So What Can We Do About It? An Actionable Survey. arXiv:2210.07700, 2022. \url{https://arxiv.org/abs/2210.07700}
\bibitem{ref337} MBIAS: Mitigating Bias in Large Language Models While Retaining Context. arXiv:2405.11290, 2024. \url{https://arxiv.org/abs/2405.11290}
\bibitem{ref338} A Holistic Approach to Undesired Content Detection in the Real World. arXiv:2208.03274, 2022. \url{https://arxiv.org/abs/2208.03274}
\bibitem{ref339} Characterizing and Evaluating the Reliability of LLMs against Jailbreak Attacks. arXiv:2408.09326, 2024. \url{https://arxiv.org/abs/2408.09326}
\bibitem{ref340} Detectors for Safe and Reliable LLMs: Implementations, Uses, and Limitations. arXiv:2403.06009, 2024. \url{https://arxiv.org/abs/2403.06009}
\bibitem{ref341} Risk Taxonomy, Mitigation, and Assessment Benchmarks of Large Language Model Systems. arXiv:2401.05778, 2024. \url{https://arxiv.org/abs/2401.05778}
\bibitem{ref342} LaMPP: Language Models as Probabilistic Priors for Perception and Action. arXiv:2302.02801, 2023. \url{https://arxiv.org/abs/2302.02801}
\bibitem{ref343} NLX-GPT: A Model for Natural Language Explanations in Vision and Vision-Language Tasks. arXiv:2203.05081, 2022. \url{https://arxiv.org/abs/2203.05081}
\bibitem{ref344} T-Eval Evaluating the Tool Utilization Capability of Large Language Models Step by Step. arXiv:2312.14033, 2023. \url{https://arxiv.org/abs/2312.14033}
\bibitem{ref345} Sentiment Reasoning for Healthcare. arXiv:2407.21054, 2024. \url{https://arxiv.org/abs/2407.21054}
\bibitem{ref346} Investigating Layer Importance in Large Language Models. arXiv:2409.14381, 2024. \url{https://arxiv.org/abs/2409.14381}
\bibitem{ref347} Holistic Evaluation of Language Models. arXiv:2211.09110, 2022. \url{https://arxiv.org/abs/2211.09110}
\bibitem{ref348} Probing LLMs for Joint Encoding of Linguistic Categories. arXiv:2310.18696, 2023. \url{https://arxiv.org/abs/2310.18696}
\bibitem{ref349} Faithfulness vs. Plausibility: On the (Un)Reliability of Explanations from Large Language Models. arXiv:2402.04614, 2024. \url{https://arxiv.org/abs/2402.04614}
\bibitem{ref350} Uncertainty as a Form of Transparency: Measuring, Communicating, and Using Uncertainty. doi:10.1145/3461702.3462571, 2021. \url{https://doi.org/10.1145/3461702.3462571}
\bibitem{ref351} LLM Inference Serving: Survey of Recent Advances and Opportunities. doi:10.1109/hpec62836.2024.10938426, 2024. \url{https://doi.org/10.1109/hpec62836.2024.10938426}
\bibitem{ref352} Mastering the Craft of Data Synthesis for CodeLLMs. arXiv:2411.00005, 2024. \url{https://arxiv.org/abs/2411.00005}
\bibitem{ref353} Towards Sustainable NLP: Insights from Benchmarking Inference Energy in Large Language Models. arXiv:2502.05610, 2025. \url{https://arxiv.org/abs/2502.05610}
\bibitem{ref354} The Efficiency Spectrum of Large Language Models: An Algorithmic Survey. arXiv:2312.00678, 2023. \url{https://arxiv.org/abs/2312.00678}
\bibitem{ref355} Emergent autonomous scientific research capabilities of large language models. arXiv:2304.05332, 2023. \url{https://arxiv.org/abs/2304.05332}
\bibitem{ref356} LLMs are Also Effective Embedding Models: An In-depth Overview. arXiv:2412.12591, 2024. \url{https://arxiv.org/abs/2412.12591}
\bibitem{ref357} MetaReflection: Learning Instructions for Language Agents using Past Reflections. arXiv:2405.13009, 2024. \url{https://arxiv.org/abs/2405.13009}
\bibitem{ref358} Robust Prompt Optimization for Large Language Models Against Distribution Shifts. doi:10.18653/v1/2023.emnlp-main.95, 2023. \url{https://doi.org/10.18653/v1/2023.emnlp-main.95}
\bibitem{ref359} Large Language Model Applied in Multi-agent SystemA Survey. doi:10.54254/2755-2721/109/20241330, 2024. \url{https://doi.org/10.54254/2755-2721/109/20241330}
\bibitem{ref360} Decoding Large-Language Models: A Systematic Overview of Socio-Technical Impacts, Constraints, and Emerging Questions. arXiv:2409.16974, 2024. \url{https://arxiv.org/abs/2409.16974}
\bibitem{ref361} Beyond Efficiency A Systematic Survey of Resource-Efficient Large Language Models. arXiv:2401.00625, 2024. \url{https://arxiv.org/abs/2401.00625}
\bibitem{ref362} Unveiling the Lexical Sensitivity of LLMs: Combinatorial Optimization for Prompt Enhancement. arXiv:2405.20701, 2024. \url{https://arxiv.org/abs/2405.20701}
\bibitem{ref363} Dynamic LLM-Agent Network An LLM-agent Collaboration Framework with Agent Team Optimization. arXiv:2310.02170, 2023. \url{https://arxiv.org/abs/2310.02170}
\bibitem{ref364} Agent-Pro: Learning to Evolve via Policy-Level Reflection and Optimization. arXiv:2402.17574, 2024. \url{https://arxiv.org/abs/2402.17574}
\bibitem{ref365} MPO: Boosting LLM Agents with Meta Plan Optimization. arXiv:2503.02682, 2025. \url{https://arxiv.org/abs/2503.02682}
\bibitem{ref366} OptiMUS: Scalable Optimization Modeling with (MI)LP Solvers and Large Language Models. arXiv:2402.10172, 2024. \url{https://arxiv.org/abs/2402.10172}
\bibitem{ref367} Survival of the Safest: Towards Secure Prompt Optimization through Interleaved Multi-Objective Evolution. arXiv:2410.09652, 2024. \url{https://arxiv.org/abs/2410.09652}
\bibitem{ref368} Achieving Peak Performance for Large Language Models: A Systematic Review. doi:10.1109/access.2024.3424945, 2024. \url{https://doi.org/10.1109/access.2024.3424945}
\bibitem{ref369} COSMOS: A Hybrid Adaptive Optimizer for Memory-Efficient Training of LLMs. arXiv:2502.17410, 2025. \url{https://arxiv.org/abs/2502.17410}
\bibitem{ref370} Flexora: Flexible Low Rank Adaptation for Large Language Models. arXiv:2408.10774, 2024. \url{https://arxiv.org/abs/2408.10774}
\bibitem{ref371} OptiChat: Bridging Optimization Models and Practitioners with Large Language Models. arXiv:2501.08406, 2025. \url{https://arxiv.org/abs/2501.08406}
\bibitem{ref372} Language Agents as Optimizable Graphs. arXiv:2402.16823, 2024. \url{https://arxiv.org/abs/2402.16823}
\end{thebibliography}

\end{document}
